\documentclass[11pt,a4paper]{article}
\usepackage{DDTA_METHODOLOGY_GUIDE_STYLE_R1}
\usepackage{amsmath,amssymb}

\hypersetup{
  pdftitle={DDTA Base Analysis Guide - Rebuild R7 - condition integration},
  pdfauthor={Documentation-Driven Threat Analysis research project}
}

\newcommand{\BA}{\textit{Base Analysis}}
\newcommand{\BAR}{\texttt{BAReferent}}
\newcommand{\BAP}{\texttt{BAProposition}}
\newcommand{\CurrentTag}{\ddtatag{DDTAGreen}{CURRENT / PRESERVED}}
\newcommand{\CandidateTag}{\ddtatag{DDTAAmber}{CANDIDATE / NOT ADMITTED}}
\newcommand{\OpenTag}{\ddtatag{DDTARed}{OPEN}}
\newcommand{\RejectedTag}{\ddtatag{DDTARed}{REJECTED}}
\newcommand{\RebuildTag}{\ddtatag{DDTABlue}{REBUILD CONTRACT}}
\newcommand{\PossibleStructuralRedundancy}{\texttt{POSSIBLE\_UPSTREAM\_STRUCTURAL\_REDUNDANCY}}

\newcommand{\TemporaryRebuildBanner}{%
\begin{ddtaanti}[title={APPENDICE TEMPORANEA DI REBUILD - DA ELIMINARE PRIMA DELLA GUIDA DEFINITIVA}]
Il contenuto di questa appendice governa \emph{come stiamo ricostruendo la guida}, non il funzionamento della Base Analysis. Deve restare disponibile durante il rebuild per garantire coerenza, genealogia e regression, ma va rimosso dalla guida definitiva dopo il consolidamento. Se una sua regola deve sopravvivere come metodologia BA, deve essere trasferita esplicitamente nella parte metodologica destinata alla guida definitiva prima della rimozione.
\end{ddtaanti}
}

% -----------------------------------------------------------------------------
% PAGE-INTEGRITY HASHES
% -----------------------------------------------------------------------------
\newcommand{\PageMDfive}[1]{\csname PageMDfive#1\endcsname}
%<PAGE-HASH-DEFS>
\expandafter\def\csname PageMDfive1\endcsname{0c456e68bcd056124768219bacdcab68}
\expandafter\def\csname PageMDfive2\endcsname{fa01b7db301f77cd9988907daf2b3397}
\expandafter\def\csname PageMDfive3\endcsname{3b620cf21e74f64e397656ec9f3038da}
\expandafter\def\csname PageMDfive4\endcsname{b12d8f94dc04c833c0d916453255f9be}
\expandafter\def\csname PageMDfive5\endcsname{44c236d6809ab29df535fe29ff00352d}
\expandafter\def\csname PageMDfive6\endcsname{b9514c3f03bdb13c46c35ba6ae647130}
\expandafter\def\csname PageMDfive7\endcsname{db02dca9f3877577832e39efb2e6b7c1}
\expandafter\def\csname PageMDfive8\endcsname{b4b3bc661001ccb16d9acdef681806bd}
\expandafter\def\csname PageMDfive9\endcsname{dcd3d2d1eb4eb28971b488dcd1c7cf01}
\expandafter\def\csname PageMDfive10\endcsname{9fdf9a0139ea36eaa577008c81fb109f}
\expandafter\def\csname PageMDfive11\endcsname{f0c4184f4f1ebcb7ce1ab9394c5c5209}
\expandafter\def\csname PageMDfive12\endcsname{e3200807fcbf2aa290f0748b9c099b69}
\expandafter\def\csname PageMDfive13\endcsname{b9028ef062953f4b7ab73bfb7ac7975c}
\expandafter\def\csname PageMDfive14\endcsname{849a9ae9b63b6f5b84acbf3c1015a920}
\expandafter\def\csname PageMDfive15\endcsname{c6fca40abd8cf6a72634f4b1d07fa450}
\expandafter\def\csname PageMDfive16\endcsname{83b4a5ba93a8eb218e7a0816bb7e023a}
\expandafter\def\csname PageMDfive17\endcsname{8063825cd44fdd393ef91bf9a5957e86}
\expandafter\def\csname PageMDfive18\endcsname{f23c4cd3c7776e49c0d794a9ecfef9a9}
\expandafter\def\csname PageMDfive19\endcsname{bc3e4197b5ba6fc299035dd3ee55f0fa}
\expandafter\def\csname PageMDfive20\endcsname{90bad45a1ddf4546a84ffcaafe19514f}
\expandafter\def\csname PageMDfive21\endcsname{a4a5a63a56b233dc9685ef8ade259ac3}
\expandafter\def\csname PageMDfive22\endcsname{0b13ca345d9009ea1359e515bbbdf161}
\expandafter\def\csname PageMDfive23\endcsname{e33af8b0be91d76075d6f8d588152778}
\expandafter\def\csname PageMDfive24\endcsname{e3dfd0ae1ea4b8fc206e94586671531d}
\expandafter\def\csname PageMDfive25\endcsname{7fe9e1357a2287afc0bcd8cdc5d05cad}
\expandafter\def\csname PageMDfive26\endcsname{ca23e6876cbd5a9485a716d5b3e2a2a7}
\expandafter\def\csname PageMDfive27\endcsname{5b77098bebe4aee19463db63e451b2de}
\expandafter\def\csname PageMDfive28\endcsname{f59be85e2e26e94598d971d526ab5349}
\expandafter\def\csname PageMDfive29\endcsname{daa47270d6924161b6041ddc1f28526b}
\expandafter\def\csname PageMDfive30\endcsname{fbf9ff38b0f615ee2066e351a264262d}
\expandafter\def\csname PageMDfive31\endcsname{124230f8d8f7900c16fc39a88a25fa5c}
\expandafter\def\csname PageMDfive32\endcsname{a04b76a31302078a1ca5e5411d643f1e}
\expandafter\def\csname PageMDfive33\endcsname{a3497f18dc46cda6fa6fc5caf9ce25b3}
\expandafter\def\csname PageMDfive34\endcsname{90363cec59b9cd817ff464809b0071e5}
\expandafter\def\csname PageMDfive35\endcsname{33bf51c48bbe73a21ea78aaa81bce926}
\expandafter\def\csname PageMDfive36\endcsname{cc6deb21f9c4c0bc45bf043617c58d0c}
\expandafter\def\csname PageMDfive37\endcsname{929ad404323716bdb65630c172f8b864}
\expandafter\def\csname PageMDfive38\endcsname{29011b2ce81135be54d4b0002b9c69c8}
\expandafter\def\csname PageMDfive39\endcsname{7abd09f421a17323e61078e606c76cc5}
\expandafter\def\csname PageMDfive40\endcsname{b0e191c6045d2c9e9e1f7237f8a64722}
\expandafter\def\csname PageMDfive41\endcsname{fd686d1106dfb09f995cbd206d72f204}
\expandafter\def\csname PageMDfive42\endcsname{7bf58985b2b7e03d4fa8cedf3a820e17}
\expandafter\def\csname PageMDfive43\endcsname{376f50ac864805c54ebb05510ddd2437}
\expandafter\def\csname PageMDfive44\endcsname{787e4c5ce9931b03135f6fe32337ab0c}
\expandafter\def\csname PageMDfive45\endcsname{878f6fa9ebb37d3e1c92b53232d28d37}
\expandafter\def\csname PageMDfive46\endcsname{1d9565d3ea82d81047cb85703c2209f8}
\expandafter\def\csname PageMDfive47\endcsname{541905031f0138d136150e7ae4a8c356}
\expandafter\def\csname PageMDfive48\endcsname{452922a9ba3eda4dd4fa6a64e3e8f262}
\expandafter\def\csname PageMDfive49\endcsname{ea668a6b4e75d6ae2af23bddd09f313e}
\expandafter\def\csname PageMDfive50\endcsname{8eaecceab767d383792c89556a2de0a5}
\expandafter\def\csname PageMDfive51\endcsname{19e590e985946f03fa9d509114958c3b}
\expandafter\def\csname PageMDfive52\endcsname{c1b9d6fafafa9462c73f4489fad3664d}
\expandafter\def\csname PageMDfive53\endcsname{ed20ede9432d24c5fee8c50f3d80f862}
\expandafter\def\csname PageMDfive54\endcsname{9866d0d6f214e0b181ad9015f5137ad9}
\expandafter\def\csname PageMDfive55\endcsname{52f8b97d73181ebc63f4ffc2419beed0}
\expandafter\def\csname PageMDfive56\endcsname{1e66cad562b25b4881d43e62c217878b}
\expandafter\def\csname PageMDfive57\endcsname{b4c3dc1c60fb207da9be278f13f0263e}
\expandafter\def\csname PageMDfive58\endcsname{9ac0774d88929da648825995c1076378}
\expandafter\def\csname PageMDfive59\endcsname{ff13530c402da52d442b404aa8b443d7}
\expandafter\def\csname PageMDfive60\endcsname{d5d06eb1db7d0bce9416648ab93e2c02}
\expandafter\def\csname PageMDfive61\endcsname{a6de99c516eda39f3ac29f1a6246bb90}
\expandafter\def\csname PageMDfive62\endcsname{ec89bcdef246536ca693660807fae78d}
\expandafter\def\csname PageMDfive63\endcsname{5693f684c29fcc1a655605d77b85877d}
\expandafter\def\csname PageMDfive64\endcsname{3cef93fdcd6726629e1dd667c5c0f775}
\expandafter\def\csname PageMDfive65\endcsname{7956aba2ac06ebcc3c7973d91b70469a}
\expandafter\def\csname PageMDfive66\endcsname{4d65df60cd00baef5ce9044dbbc85f99}
\expandafter\def\csname PageMDfive67\endcsname{ecda2667c8851cd4a8a88c0b5ab7a817}
\expandafter\def\csname PageMDfive68\endcsname{838b86747edfaf9255eba656359ba686}
\expandafter\def\csname PageMDfive69\endcsname{ae76776a0123bf292f75f6e6c8e85873}
\expandafter\def\csname PageMDfive70\endcsname{81fa8e48c04924374f844530857a7961}
\expandafter\def\csname PageMDfive71\endcsname{a0be33fafc8818cbc8aa9a228800ac06}
\expandafter\def\csname PageMDfive72\endcsname{78df03da071146eaf0784e3cb663f8e1}
\expandafter\def\csname PageMDfive73\endcsname{e71b51c4b33009106e6b6a6244365314}
\expandafter\def\csname PageMDfive74\endcsname{31a35f3c5c6d7bc55d535686970e2869}
\expandafter\def\csname PageMDfive75\endcsname{f248e4c48942111fde0c866c423d5e24}
\expandafter\def\csname PageMDfive76\endcsname{ba2f8053b6612bc2a52396c2fd3ec949}
\expandafter\def\csname PageMDfive77\endcsname{3915eaee02f714ac47e0dfa764b2836f}
\expandafter\def\csname PageMDfive78\endcsname{5475a02c3223c3fac62a7eebe62f6be2}
\expandafter\def\csname PageMDfive79\endcsname{e1b1becfc138e20474529fae47de6a05}
\expandafter\def\csname PageMDfive80\endcsname{435e9ecddb23d230df9cc32766e10610}
\expandafter\def\csname PageMDfive81\endcsname{9d02d27cc2ad8bb4e1b30d441667e854}
\expandafter\def\csname PageMDfive82\endcsname{90b6659e5a0105e0cd7268d3f76137f6}
\expandafter\def\csname PageMDfive83\endcsname{4076e2f7a00855a7ac650ff3ac7ab044}
\expandafter\def\csname PageMDfive84\endcsname{080f4ff9e7080b3504b21842c5a514db}
\expandafter\def\csname PageMDfive85\endcsname{13cdd31df41e1d75e405e6ba020261fa}
%</PAGE-HASH-DEFS>

\newcommand{\PageHashFooter}[1]{%
  \fancyfoot[C]{}%
  \fancyfoot[R]{\sffamily\scriptsize\color{DDTAGray}Pag. \thepage\quad MD5 contenuto: \texttt{#1}}%
}
\newcommand{\IntegrityRow}[3]{#1 & #2 & {\footnotesize\ttfamily #3} \\
}

\begin{document}
\selectlanguage{italian}
\fancyhead[L]{\sffamily\small\color{DDTAGray}DDTA Base Analysis Guide}
\fancyhead[R]{\sffamily\small\color{DDTABlue}Rebuild R7 - CANDIDATE}
\fancyfoot[L]{\sffamily\scriptsize\color{DDTAGray}Cumulative rebuild - human-readable method, controlled semantics}

%<PAGE-DESC:1>Frontespizio della Base Analysis Guide Rebuild R7
%<PAGE-CONTENT:1>
\begin{titlepage}
\thispagestyle{empty}
\vspace*{12mm}
{\sffamily\bfseries\color{DDTANavy}\fontsize{15}{18}\selectfont Documentation-Driven Threat Analysis}\par
\vspace{5mm}
{\sffamily\bfseries\color{DDTANavy}\fontsize{28}{32}\selectfont Base Analysis Guide}\par
\vspace{2mm}
{\sffamily\color{DDTABlue}\fontsize{18}{22}\selectfont Rebuild R7 - Candidate cumulative reconstruction}\par

\vspace{14mm}
\begin{tcolorbox}[enhanced,arc=3pt,boxrule=0pt,colback=DDTANavy,left=13pt,right=13pt,top=12pt,bottom=12pt]
{\color{white}\sffamily\bfseries\large Ricostruzione progressiva della guida Base Analysis}\par
\vspace{2mm}
{\color{white!88}\small Ogni costrutto viene ricostruito a partire dall'authority BA corrente e dalle guide precedenti, spiegato in forma generale e soltanto dopo verificato contro esempi aggiornati. Il caso DermaTriage e' evidence applicativa e worked example: non e' la fonte da cui nasce la semantica del costrutto.}
\end{tcolorbox}

\vspace{5mm}
\begin{ddtaopen}[title={RIQUADRO TEMPORANEO - WORKING SET E RIFERIMENTI - DA CANCELLARE PRIMA DELLA GUIDA DEFINITIVA}]
{\scriptsize
\textbf{Attivi:} \nolinkurl{DDTA_DOCUMENTATION_AUTHORING_GUIDE_R7_REBUILD_R8.tex}; \nolinkurl{DDTA_BASE_ANALYSIS_GUIDE_REBUILD_R7.tex}; \nolinkurl{DDTA_DERMATRIAGE_PARALLEL_CASE_STUDY_R25_FR03_SPLIT_REVIEW_R1.tex}.\par
\textbf{Piano corrente:} \nolinkurl{DDTA_R25_BASE_ANALYSIS_GUIDE_REBUILD_WORK_PLAN_R5.md}; working analysis: \nolinkurl{DDTA_R25_DERMATRIAGE_CASE_STUDY_WORKING_ANALYSIS_R1.md}.\par
\textbf{Riferimenti metodologici:} Operational Guide R3 + BA0--BA5 come authority corrente; Core/Complete BA Candidate R1 e precedenti guide operative come research evidence per la ricostruzione didattica.\par
\textbf{Fase corrente:} mantenere R7 cumulativa e ricostruire soltanto i costrutti richiesti dalla BA corrente. \texttt{transfer}, \texttt{produce}, \texttt{realize}, \texttt{constrain} e lo scoped modifier \texttt{condition} restano ricostruiti; questo incremento testa, come candidate non ammessa, la rilocazione dello storico \texttt{decisionRule} dentro \texttt{produce} come \texttt{operatorStructure.decisionRule}, usando DermaTriage R25 e il regression case Facial Access come pressure.
}
\end{ddtaopen}

\vfill
\begin{tabularx}{0.93\textwidth}{L{46mm}Y}
\toprule
\textbf{Stato} & \texttt{CANDIDATE / NON-NORMATIVE / CUMULATIVE REBUILD} \\
\textbf{Repository baseline} & \texttt{b15990088769e13c7f776a9591b9f0c1d5e11053} \\
\textbf{Predecessore cumulativo} & \nolinkurl{DDTA_BASE_ANALYSIS_GUIDE_REBUILD_R6} \\
\textbf{Authority BA corrente} & R3 Operational Guide + BA0 R1 / BA1 R1 / BA2 R3 / BA3 R1 / BA4 R1 / BA5 R1 \\
\textbf{Scopo di R7} & Ricostruzione cumulativa dei costrutti BA richiesti dal caso. I quattro operatori ricostruiti e \texttt{condition} restano invariati; questo incremento testa la candidate \texttt{operatorStructure.decisionRule}, cioe' la rilocazione dello stesso costrutto storico \texttt{decisionRule} dentro \texttt{produce} senza duplicare actor/input/result. \\
\textbf{Authority change} & Nessuno. \texttt{decisionRule} resta authority corrente come \texttt{semanticOperatorKey} top-level; la sua rilocazione candidate come \texttt{operatorStructure.decisionRule} richiede regression + checkpoint prima di promotion. \\
\bottomrule
\end{tabularx}

\vfill
{\sffamily\scriptsize\color{DDTAGray}MD5 contenuto pagina 1: \texttt{\PageMDfive{1}}}\par
\end{titlepage}
%</PAGE-CONTENT:1>

\clearpage

%<PAGE-DESC:2>1 - Che cos'e' la Base Analysis
%<PAGE-CONTENT:2>
\PageHashFooter{\PageMDfive{2}}
\section{Che cos'e' la Base Analysis}

La \BA{} e' una rappresentazione analitica condivisa del significato di progetto. Il suo compito e' rendere esplicite, in una forma comune, le informazioni che devono poter essere riutilizzate, interrogate, confrontate o proiettate verso analisi e viste successive.

Il punto di partenza resta sempre la documentazione. E' la documentazione a stabilire che cosa il progetto governa; la Base Analysis organizza quel significato senza sostituirsi alla sua authority. In questo modo consumer diversi possono lavorare sulla stessa base semantica senza dover reinterpretare ogni volta la prosa sorgente.

\begin{ddtaprinciple}[title={Idea fondamentale}]
La BA rende analiticamente esplicito un significato gia' governato. Non crea project truth e non diventa una fonte alternativa alla documentazione.
\end{ddtaprinciple}

\subsection{Che cosa resta fuori dalla BA}

Una Base Analysis completa non coincide con una copia strutturata di tutto cio' che compare nei documenti. Diagrammi architetturali, DFD, modelli STRIDE e viste di progetto possono essere prodotti a valle, ma non definiscono da soli che cosa debba entrare nella BA. Allo stesso modo, la semplice presenza di un sostantivo, di un verbo, di una tecnologia o di un valore nel testo non basta a giustificare un nuovo elemento analitico.

Il criterio e' semantico: rappresentiamo cio' che deve restare distinguibile e riusabile per gli scopi dichiarati, mantenendo il livello di dettaglio effettivamente governato dal progetto.
%</PAGE-CONTENT:2>

\clearpage

%<PAGE-DESC:3>1 - Perche' esiste e principi di base
%<PAGE-CONTENT:3>
\PageHashFooter{\PageMDfive{3}}
\subsection{Perche' esiste}

Uno stesso significato di progetto puo' servire a consumer molto diversi. Una vista del sistema puo' aver bisogno di mostrare responsabilita' e interazioni; la traceability deve poter risalire alle fonti; il change impact deve riconoscere che cosa viene toccato da una modifica; test reasoning e threat analysis possono porre domande ancora differenti.

Se ciascun consumer ripartisse direttamente dalla prosa, aumenterebbe il rischio di interpretazioni incompatibili. La BA introduce quindi una base analitica comune: il significato viene ricostruito una volta in modo controllato e poi riutilizzato nelle diverse attivita' downstream.

\subsection{Authority, minimalita' e feedback}

Tre esigenze devono rimanere in equilibrio. La prima e' preservare l'authority documentale: cio' che compare nella BA deve poter essere ricondotto al significato governato. La seconda e' la minimalita': aggiungere struttura che non preserva alcuna distinzione utile rende la rappresentazione piu' pesante senza renderla piu' corretta. La terza e' il feedback: quando una vista o un'analisi incontra un problema, la difficolta' deve poter essere riportata al livello che puo' davvero risolverla.

Questi tre principi spiegano perche' la BA e' contemporaneamente una rappresentazione condivisa e un confine: abbastanza ricca da essere riusata, ma non tanto libera da trasformarsi in una seconda documentazione di progetto.
%</PAGE-CONTENT:3>

\clearpage

%<PAGE-DESC:4>1 - Direzione dell'authority e feedback
%<PAGE-CONTENT:4>
\PageHashFooter{\PageMDfive{4}}
\subsection{Direzione dell'authority}

La governed project documentation rimane la project authority. Una Base Analysis accettata costituisce invece la shared analytical baseline: una rappresentazione comune che puo' essere consumata da viste, projection, threat method, test model e altri strumenti analitici.

La direzione dell'authority non cambia quando un consumer downstream scopre una difficolta'. Quel consumer puo' mostrare che una relazione e' ambigua, che manca un'informazione necessaria o che la BA non riesce a preservare una distinzione. Il problema torna allora alla BA come domanda di review. Se la BA dispone gia' del significato necessario, si corregge la rappresentazione; se il significato non e' governato, la domanda viene riportata alla documentazione.

\begin{ddtacore}[title={Feedback governato}]
Una difficolta' downstream puo' provocare una review della BA e, quando necessario, una domanda documentale. Solo una modifica governata della documentazione puo' cambiare il project meaning; dopo tale modifica la BA viene rivalidata e le projection interessate vengono ricostruite.
\end{ddtacore}

Il feedback migliora quindi la qualita' del modello e della documentazione senza invertire la sorgente dell'authority.
%</PAGE-CONTENT:4>

\clearpage

%<PAGE-DESC:5>1 - Minimum sufficient BA, NOT SPECIFIED e gap
%<PAGE-CONTENT:5>
\PageHashFooter{\PageMDfive{5}}
\subsection{Minimum sufficient BA}

Una BA piu' dettagliata non e' automaticamente una BA migliore. Il livello giusto e' quello che rende esplicite le distinzioni che devono davvero sopravvivere all'analisi.

Vale la pena introdurre un significato quando serve a mantenere un'identita' riutilizzabile, a distinguere una relazione o una condizione governata, a sostenere traceability e change/revalidation oppure a rispondere a una domanda downstream materiale usando informazioni gia' presenti nella documentazione. Se nessuna di queste esigenze esiste, il dettaglio puo' restare fuori dalla rappresentazione condivisa.

Una domanda downstream puo' quindi motivare l'esplicitazione di un significato gia' governato; non puo' creare il significato che le serve.

\begin{ddtacheck}[title={Stopping rule iniziale}]
Se togliendo un dettaglio non perdiamo una distinzione governata, un'identita' da riusare, una relazione necessaria o la possibilita' di rispondere a una domanda materiale nello scope dichiarato, quel dettaglio non viene aggiunto soltanto per rendere il modello piu' completo graficamente.
\end{ddtacheck}

\begin{ddtaopen}[title={\texttt{NOT SPECIFIED} e documentation gap}]
Quando la rappresentazione richiederebbe un fatto che la documentazione non governa, il risultato corretto puo' essere \texttt{NOT SPECIFIED} oppure una diagnosis. \texttt{NOT SPECIFIED} non segnala automaticamente un difetto: quel fatto potrebbe semplicemente non appartenere allo scope del progetto. Diventa una documentation question solo quando l'assenza impedisce di esprimere un significato che il progetto dovrebbe governare.
\end{ddtaopen}
%</PAGE-CONTENT:5>

\clearpage

%<PAGE-DESC:6>2 - BA View / Projection Contract e source types
%<PAGE-CONTENT:6>
\PageHashFooter{\PageMDfive{6}}
\section{BA View / Projection Contract}

Una stessa Base Analysis deve poter alimentare viste differenti senza trasformare ogni vista in una nuova fonte di significato. Per questo introduciamo un \textbf{View Contract}: uno strato esplicito che dichiara quale parte del significato disponibile viene selezionata, come viene proiettata e con quale notazione viene resa visibile.

Quando la sorgente e' la BA, il percorso puo' essere riassunto con un piccolo schema:

\begin{center}
\begin{tcolorbox}[width=0.84\textwidth,colback=DDTALightGray,colframe=DDTAGray!55,boxrule=0.5pt,arc=2pt]
\centering\sffamily
Governed Documentation \(\longrightarrow\) Accepted Base Analysis \(\longrightarrow\) View Contract \(\longrightarrow\) Rendered View
\end{tcolorbox}
\end{center}

Non tutte le viste, pero', partono esclusivamente dalla BA. Una vista della gerarchia documentale puo' leggere direttamente MR, Decision e FunctionalRequirement; una vista di traceability puo' invece combinare documentazione e rappresentazione analitica. Il contract deve rendere esplicita questa scelta.

\subsection{Da quale sorgente parte la view}

Una \texttt{BA\_ONLY} view usa soltanto semantica presente nella Base Analysis. Una \texttt{DOCUMENTATION\_ONLY} view rappresenta la struttura o il contenuto governato della documentazione senza duplicarlo artificialmente nella BA. Una \texttt{DOCUMENTATION\_PLUS\_BA\_TRACEABILITY} view mette infine in relazione gli elementi documentali con gli elementi analitici che ne rappresentano il significato.

Questa distinzione evita di inserire nella BA informazioni che servono soltanto a rendere possibile una particolare visualizzazione.
%</PAGE-CONTENT:6>

\clearpage

%<PAGE-DESC:7>2 - Mermaid e campi minimi della view
%<PAGE-CONTENT:7>
\PageHashFooter{\PageMDfive{7}}
\subsection{Mermaid come primo target}

Mermaid viene adottato come primo target di rendering testuale da investigare. E' utile perche' mantiene i diagrammi come artifact versionabili nel repository e offre piu' famiglie di rappresentazione senza imporre un editor grafico separato.

Mermaid resta pero' una notazione di output. Se la sua sintassi rende comodo introdurre un nodo, una relazione, un participant o un subgraph, quella comodita' non basta a creare nuova semantica nella BA. Prima viene il significato; poi il View Contract decide come renderlo.

\begin{ddtacore}[title={Semantic determinism e no reverse inference}]
A parita' di Base Analysis accettata e di projection profile, una view canonica MUST esporre lo stesso inventario semantico: stessi referent inclusi, stesse proposition, stessi role/edge meaning, stessi modifier/guard e stessi grouping semanticamente richiesti. Il renderer puo' spostare geometricamente gli oggetti gia' collegati senza violare questo requisito. Una scelta grafica non puo' aggiungere semantica ne' diventare la ragione per modificare la BA.
\end{ddtacore}

\subsection{Che cosa deve dichiarare una view}

Il contract deve identificare la view tramite \texttt{viewId} e \texttt{viewType}; deve poi indicare la sorgente con \texttt{sourceType} e \texttt{sourceBaseline} e delimitare lo \texttt{scope}. Quando la sorgente comprende la BA, specifica quali \texttt{includedReferents} e \texttt{includedPropositionsOrLocalStructures} vengono usati. Le trasformazioni sono rese esplicite nelle \texttt{projectionRules}; \texttt{coverageMode} e \texttt{omittedSemantics} spiegano quanto della sorgente viene mostrato, mentre \texttt{renderer} e \texttt{renderedArtifact} identificano la forma di output.

I campi che non hanno senso per uno specifico \texttt{sourceType} vengono marcati come non applicabili, invece di essere riempiti artificialmente.

Una struttura locale presente nella BA, inclusa una futura \texttt{operatorStructure}, non deve essere trasformata in uno stato \emph{collapsed} del modello per semplificare il diagramma. Il projection profile puo' ometterne completamente il rendering quando non e' pertinente alla view; \texttt{coverageMode}/\texttt{omittedSemantics} rendono esplicita l'omissione senza modificare la BA sottostante.
%</PAGE-CONTENT:7>

\clearpage

%<PAGE-DESC:8>2 - Coverage mode e famiglie di viste
%<PAGE-CONTENT:8>
\PageHashFooter{\PageMDfive{8}}
\subsection{Coverage mode}

BA4 distingue due modi di dichiarare la copertura. Una view \texttt{EXHAUSTIVE\_FOR\_DECLARED\_SCOPE} afferma di rappresentare tutto il significato che appartiene allo scope definito ed e' compatibile con quel tipo di vista. Una view \texttt{SELECTIVE} mostra deliberatamente soltanto una parte del significato disponibile.

La seconda forma e' normale e spesso desiderabile: una vista di responsabilita', per esempio, non deve necessariamente mostrare tutte le informazioni, gli stati o i vincoli presenti nella stessa BA. Cio' che viene omesso da una view selettiva non diventa per questo falso, assente o non governato.

\subsection{Famiglie di viste da testare}

Le prime prove dovranno coprire piu' prospettive. Alcune riguardano la struttura del progetto, come capability, responsabilita' e dipendenze. Altre mostrano aspetti dinamici o relazionali, per esempio information flow, service interaction, state/lifecycle e, quando la semantica disponibile lo consente, pipeline o processi.

Un secondo gruppo riguarda invece la documentazione: gerarchia documentale, traceability fra documentazione e BA, copertura e gap. Non fissiamo ancora il mapping di queste viste; la loro utilita' verra' verificata mentre costruiamo la BA di DermaTriage.

Il fatto che una famiglia di viste sia utile non dimostra che la BA corrente contenga gia' tutto cio' che serve a produrla.
%</PAGE-CONTENT:8>

\clearpage

%<PAGE-DESC:9>2 - View gap e mapping discipline
%<PAGE-CONTENT:9>
\PageHashFooter{\PageMDfive{9}}
\subsection{Quando una view non e' producibile}

Una view puo' mettere in evidenza un'assenza senza avere il diritto di colmarla. Se il significato richiesto non e' presente nella BA, il View Contract registra \texttt{VIEW REQUIRES MISSING BA SEMANTICS}. Se, a sua volta, la BA non puo' ricavare quel significato dalla documentazione governata, il problema diventa \texttt{BA REQUIRES MISSING GOVERNED MEANING}.

Questa distinzione permette di riportare la difficolta' al livello corretto invece di risolverla con una convenzione grafica o una supposizione del renderer.

\subsection{Come una relazione entra nella view}

Ogni elemento visuale che porta significato deve poter essere giustificato da una regola di projection dichiarata. Un edge Mermaid, per esempio, non e' valido perche' esiste una freccia con un nome simile nella notazione: deve essere il risultato di una trasformazione esplicita di semantica documentale o BA.

La notazione decide come mostrare il risultato della projection; non decide che cosa quel risultato significa. Per questo i mapping fra DDTA, BA e renderer verranno definiti e testati separatamente, invece di essere dedotti per somiglianza terminologica.
%</PAGE-CONTENT:9>

\clearpage

%<PAGE-DESC:10>3 - Dal testo documentale al significato governato
%<PAGE-CONTENT:10>
\PageHashFooter{\PageMDfive{10}}
\section{Dal testo documentale alla Base Analysis}

La Base Analysis non nasce traducendo meccanicamente le frasi della documentazione in nodi e relazioni. Prima di scegliere un elemento BA dobbiamo capire quale significato il testo stabilisce realmente e quali parti di quel significato devono rimanere distinguibili nell'analisi.

La forma linguistica e la struttura analitica non coincidono necessariamente. Una frase puo' contenere piu' fatti rilevanti; piu' frasi possono contribuire allo stesso significato; una parola tecnica puo' descrivere una realizzazione senza richiedere una nuova identita' analitica.

\begin{ddtaprinciple}[title={Prima il significato, poi il costrutto}]
La domanda iniziale non e' quale operatore somigli al verbo che stiamo leggendo. E' quale significato di progetto resti vero indipendentemente da come la frase e' stata scritta.
\end{ddtaprinciple}

Da qui si procede per gradi. Si parte dalla documentazione governata, si ricostruisce il significato che essa effettivamente stabilisce, si isolano i fatti semantici minimi che servono e si chiarisce quali identita' e quali asserzioni hanno bisogno di una forma analitica. Solo allora si materializzano BAReferent e BAProposition e, in un passaggio successivo, si sceglie il costrutto formale piu' adatto.

\begin{ddtaanti}[title={Nessun word-to-model mapping}]
La presenza di un sostantivo non crea automaticamente un BAReferent; un verbo non seleziona da solo un operatore; il confine di una frase non determina il confine di una BAProposition.
\end{ddtaanti}
%</PAGE-CONTENT:10>

\clearpage

%<PAGE-DESC:11>3 - Governed semantic fact
%<PAGE-CONTENT:11>
\PageHashFooter{\PageMDfive{11}}
\subsection{Governed semantic fact}

Chiamiamo \emph{governed semantic fact} il piu' piccolo significato governato che e' utile isolare per decidere come rappresentarlo nella BA. Non e' un nuovo tipo di elemento BA: e' un passaggio di analisi che ci permette di separare la comprensione del testo dalla scelta della sintassi analitica.

Un fatto candidato deve sempre poter essere giustificato tornando alla documentazione. Se per formularlo dobbiamo aggiungere un attore, un valore, un ordine, una causalita', una responsabilita' o una condizione non stabilita dal progetto, non stiamo piu' isolando il significato: lo stiamo completando.

\begin{ddtaexample}[title={DermaTriage - dominio della priorita'}]
La documentazione stabilisce che, quando DermaTriage produce una priorita' operativa di triage, il valore appartiene all'insieme \texttt{P1, P2, P3, P4}. Possiamo quindi isolare il fatto che \textbf{OperationalPriority ha come dominio ammesso P1--P4}. In questo momento non serve ancora decidere se \texttt{OperationalPriority} debba essere un BAReferent o quale operatore rappresentera' il vincolo.
\end{ddtaexample}

\begin{ddtaexample}[title={DermaTriage - percorso senza immagine}]
Nel percorso symptom-only emergono almeno due significati distinguibili: l'assenza dell'immagine non interrompe il triage e l'urgenza viene determinata usando le informazioni sintomatologiche disponibili sullo stesso caso. Tenerli separati evita di trasformare immediatamente l'intero paragrafo in una sola asserzione sovraccarica.
\end{ddtaexample}
%</PAGE-CONTENT:11>

\clearpage

%<PAGE-DESC:12>3 - Segmentazione, ambiguita e livello documentale
%<PAGE-CONTENT:12>
\PageHashFooter{\PageMDfive{12}}
\subsection{Come isolare i fatti senza inventarli}

La segmentazione serve a preservare differenze che potrebbero avere identita', relazioni, condizioni o provenance diverse; non serve a spezzare il testo fino a ottenere una proposition per ogni frase o per ogni verbo.

Il lavoro parte dal commitment governato e separa soltanto le distinzioni che potrebbero cambiare o essere riutilizzate indipendentemente. Le parti che perderebbero significato se isolate restano unite. Se il testo ammette piu' letture materialmente diverse, l'ambiguita' rimane visibile: non viene risolta usando conoscenza di dominio, convenzioni tecniche o dettagli dell'implementazione.

\begin{ddtacheck}[title={Delete test sul fatto candidato}]
Se elimino questa distinzione prima di costruire la BA, perdo qualcosa che la documentazione governa e che potrebbe diventare rilevante per reuse, traceability, change o una domanda downstream? Se la risposta e' no, probabilmente non serve isolarla come fatto autonomo.
\end{ddtacheck}

\subsection{Il livello documentale resta visibile}

MR, Decision e FunctionalRequirement possono governare la stessa area con livelli di precisione diversi. La BA non appiattisce questa progressione. Un MR puo' stabilire un risultato e i suoi input in termini generali; una Decision puo' restringere una strategia; un FR puo' precisare un dominio, una condizione o una regola operativa.

Quando la documentazione introduce un dettaglio piu' specifico, la BA puo' aggiungere la relativa asserzione senza cancellare automaticamente il significato piu' generale che rimane valido.
%</PAGE-CONTENT:12>

\clearpage

%<PAGE-DESC:13>4 - Elementi fondamentali della Base Analysis
%<PAGE-CONTENT:13>
\PageHashFooter{\PageMDfive{13}}
\section{Gli elementi fondamentali della Base Analysis}

La Base Analysis distingue due famiglie fondamentali di elementi: \BAR{} e \BAP{}. La prima serve a dare identita' analitica a un significato di progetto; la seconda a dare identita' a un'asserzione analitica su uno o piu' significati.

\texttt{BAE} puo' essere usato come termine ombrello per riferirsi agli elementi della Base Analysis, ma non introduce una terza famiglia.

\begin{ddtaprinciple}[title={Due domande diverse}]
Quando incontriamo un significato chiediamo se deve poter essere riconosciuto e riutilizzato come identita' analitica. Quando incontriamo un'asserzione chiediamo invece se deve poter essere tracciata, revisionata o diagnosticata indipendentemente. La prima domanda porta a un BAReferent; la seconda a una BAProposition.
\end{ddtaprinciple}

Questa distinzione e' semantica, non implementativa. Un tool puo' memorizzare referent e proposition in classi, tabelle o grafi differenti; cio' che deve restare stabile e' la differenza tra la cosa cui attribuiamo identita' e l'asserzione che facciamo su di essa.
%</PAGE-CONTENT:13>

\clearpage

%<PAGE-DESC:14>4 - BAReferent: identita analitica riusabile
%<PAGE-CONTENT:14>
\PageHashFooter{\PageMDfive{14}}
\subsection{BAReferent: dare identita' a un significato}

Un \BAR{} e' un'unita' di significato di progetto che deve poter essere riconosciuta e riutilizzata indipendentemente nell'analisi. Non e' limitato agli oggetti software o ai componenti runtime: puo' rappresentare una capability, un'informazione, un risultato, un comportamento, un evento, uno stato, un contratto, uno store, un boundary o un altro significato che abbia bisogno di identita' stabile.

\begin{ddtacheck}[title={Identity test}]
La domanda utile non e' ``questo nome compare nel testo?'', ma ``questo significato deve vivere oltre l'asserzione in cui l'ho incontrato?''. L'identita' separata diventa giustificata quando il significato deve essere riutilizzato in piu' proposition, qualificato o vincolato, correlato o referenziato, confrontato fra baseline, seguito attraverso change/revalidation oppure interrogato come oggetto distinto da un consumer downstream.
\end{ddtacheck}

Se nessuna di queste esigenze esiste e la rimozione dell'identita' non fa perdere significato governato, un BAReferent separato puo' essere superfluo.

La presenza lessicale, da sola, non basta: tecnologie, valori, campi e sostantivi diventano referent solo quando serve davvero un'identita' analitica indipendente.
%</PAGE-CONTENT:14>

\clearpage

%<PAGE-DESC:15>4 - BAReferent: esempi DermaTriage e confini
%<PAGE-CONTENT:15>
\PageHashFooter{\PageMDfive{15}}
\subsection{BAReferent: prima applicazione a DermaTriage}

La documentazione consolidata di MR-01 offre gia' alcuni buoni esempi per applicare l'identity test senza congelare prematuramente il registry finale.

\textbf{DermaTriage.} Compare ripetutamente come soggetto di obblighi e relazioni; e' quindi un forte candidato a identita' BA riusabile.

\textbf{Caso dermatologico.} Collega immagine, sintomi e risultati. La necessita' di una sua identita' autonoma dipendera' soprattutto da come dovremo preservare correlation e same-case semantics nelle proposition successive.

\textbf{Urgenza e priorita' operativa.} Sono risultati governati distinti. L'urgenza viene prodotta e riutilizzata nei passaggi successivi; la priorita' e' ulteriormente vincolata alla P-scale. Entrambe sono candidate forti a identita' riusabile, mentre resta da decidere il livello corretto fra concetto e occurrence per singolo caso.

\textbf{Il valore P1.} Il fatto che sia governato non lo trasforma automaticamente in un BAReferent. Se non deve essere riutilizzato come identita' indipendente, puo' rimanere un valore locale controllato.

\begin{ddtaopen}[title={L'esempio non decide ancora il modello finale}]
Questa prima applicazione mostra come usare l'identity test. Non congela il registry DermaTriage e non decide ancora la granularita' instance-vs-concept: saranno le proposition necessarie a mostrare quali identita' devono sopravvivere autonomamente.
\end{ddtaopen}

Differenze come actor, component, information, behavior o state non richiedono di per se' famiglie BA diverse. Finche' classificazioni, ruoli, valori e proposition riescono a preservarle onestamente, restano all'interno della stessa famiglia BAReferent.
%</PAGE-CONTENT:15>

\clearpage

%<PAGE-DESC:16>4 - BAProposition: asserzione analitica identificabile
%<PAGE-CONTENT:16>
\PageHashFooter{\PageMDfive{16}}
\subsection{BAProposition: dare identita' a una asserzione}

Una \BAP{} e' un'asserzione analitica methodology-neutral su uno o piu' BAReferent. La sua identita' permette di collegare provenance, review, diagnosis e change disposition all'asserzione stessa, senza confonderla con i significati che mette in relazione.

Questa distinzione diventa importante quando un referent rimane stabile ma cambia cio' che il progetto afferma su di esso. In quel caso non dobbiamo sostituire l'identita' del referent: possiamo introdurre, modificare, supersedere o diagnosticare le proposition interessate.

A livello BA2 la forma generale corrente prevede un \texttt{semanticOperatorKey}, una o piu' participation con roleKey scoped all'operatore e, quando ammessi, polarity, modifier e strutture locali. Le firme specifiche verranno spiegate piu' avanti costrutto per costrutto.

\begin{lstlisting}
BAProposition
    semanticOperatorKey
    participation [1..*]
        roleKey
        term
    polarity
    scopedModifier [0..*]
    operatorStructure [0..1, when admitted]
\end{lstlisting}

Il \texttt{term} di una participation e' normalmente un BAReferent. Un valore locale controllato e' sufficiente soltanto quando quel significato non richiede identita' indipendente.
%</PAGE-CONTENT:16>

\clearpage

%<PAGE-DESC:17>4 - BAReferent e BAProposition insieme
%<PAGE-CONTENT:17>
\PageHashFooter{\PageMDfive{17}}
\subsection{BAReferent e BAProposition lavorano insieme}

Prendiamo il FunctionalRequirement che governa la priorita' operativa P1--P4. Prima di formalizzare, riconosciamo due cose diverse. Da un lato esistono significati che devono poter essere riutilizzati, come DermaTriage e OperationalPriority; dall'altro la documentazione afferma qualcosa su questi significati, per esempio che una priorita' operativa puo' essere prodotta e che il suo dominio e' limitato a P1--P4.

L'identity test suggerisce quindi almeno i referent candidati \texttt{DermaTriage} e \texttt{OperationalPriority}. Le affermazioni sul risultato e sul dominio saranno invece materializzate come BAProposition quando avremo selezionato operatori, ruoli e valori compatibili con il significato governato.

\begin{ddtaexample}[title={Perche' non formalizziamo troppo presto}]
Possiamo gia' vedere che disponibilita' della priorita' e dominio ammesso sono due asserzioni distinguibili. Non serve pero' scegliere adesso \texttt{produce}, \texttt{constrain} o un altro costrutto: prima fissiamo la grammatica generale, poi riesaminiamo ogni operatore con lo stesso contratto descrittivo.
\end{ddtaexample}

Il punto da conservare e' semplice: il referent rappresenta il significato che deve continuare a esistere come identita' analitica; la proposition rappresenta cio' che affermiamo su quel significato. Un grafo o un renderer non devono confondere i due livelli soltanto per ottenere una forma visuale piu' comoda.
%</PAGE-CONTENT:17>

\clearpage

%<PAGE-DESC:18>5 - Granularita, minimalita e costruzione della BA
%<PAGE-CONTENT:18>
\PageHashFooter{\PageMDfive{18}}
\section{Granularita' e costruzione della Base Analysis}

La granularita' corretta non dipende dalla lunghezza della frase ne' dal numero di elementi che un tool riesce a disegnare. Dipende da quali distinzioni devono poter essere identificate, verificate e cambiate indipendentemente.

\subsection{Quando separare due proposition}

Due significati meritano proposition separate quando possono divergere per provenance o stato di review, per condizione di applicabilita', per impatto di un cambiamento, perche' uno dei due diventa target di una relazione successiva oppure perche' consumer downstream diversi hanno bisogno di distinguerli. Se invece le parti hanno senso soltanto insieme, la rappresentazione deve conservarne l'unita'.

Separare significa quindi preservare indipendenza semantica, non frammentare il testo.

\subsection{Procedura iniziale di costruzione}

Il lavoro puo' essere condotto in sette passaggi:
\begin{enumerate}
  \item fissare baseline, authority e scope documentale;
  \item leggere il testo source-first e ricostruire il significato governato prima di pensare alla sintassi BA;
  \item isolare soltanto i governed semantic fact che devono rimanere distinguibili;
  \item applicare l'identity test e riconoscere le asserzioni che richiedono BAReferent e BAProposition;
  \item verificare che la cardinalita' BA derivi dal significato analitico e non dalla forma dell'albero documentale; se emerge una possibile inefficienza strutturale, registrare un finding e rinviare la decisione al layer documentale;
  \item selezionare operatori, ruoli, valori e strutture locali solo dopo che identita' e asserzioni sono chiare;
  \item tornare alle fonti per verificare ogni elemento, preservare \texttt{NOT SPECIFIED} o diagnosis quando necessario e applicare infine stopping rule e forbidden-inference review.
\end{enumerate}

\begin{ddtacheck}[title={Stop prima degli operatori}]
Prima di entrare nel catalogo dei costrutti, il lettore deve saper spiegare perche' un significato richiede un BAReferent, perche' un'asserzione richiede una BAProposition e quali parti della documentazione giustificano entrambi. Se questo non e' chiaro, scegliere un operatore e' prematuro.
\end{ddtacheck}
%</PAGE-CONTENT:18>

\clearpage

%<PAGE-DESC:19>6 - Indipendenza tra struttura documentale e struttura BA
%<PAGE-CONTENT:19>
\PageHashFooter{\PageMDfive{19}}
\section{Struttura documentale, cardinalita' BA e feedback verso monte}

La documentazione DDTA organizza responsabilita', decisioni e obblighi in una struttura utile per authoring, implementazione, verifica e change management. La Base Analysis ha uno scopo diverso: rendere espliciti i significati analitici che devono essere riusati e interrogati. Per questo le due strutture non hanno una cardinalita' obbligatoriamente coincidente.

\begin{ddtaprinciple}[title={Regola di indipendenza strutturale}]
Gli elementi documentali sono sorgenti di significato governato, non un template da tradurre uno-a-uno in elementi BA.
\end{ddtaprinciple}

\begin{lstlisting}
numero di branch documentali
        !=
numero di BAReferent
        !=
numero di BAProposition
\end{lstlisting}

Piu' elementi documentali possono sostenere uno stesso significato analitico; un singolo elemento documentale puo' invece contenere piu' asserzioni che devono restare distinguibili in BA. La cardinalita' corretta emerge dai test di identita', asserzione, provenance, change e consumo downstream, non dal numero di nodi presenti nell'albero DDTA.

\begin{ddtaanti}[title={Traduzione meccanica vietata}]
Non creare un BAReferent, una BAProposition o un futuro ramo analitico soltanto perche' esiste un elemento documentale distinto. Allo stesso modo, non comprimere due significati BA distinti soltanto perche' sono scritti nello stesso requisito.
\end{ddtaanti}
%</PAGE-CONTENT:19>

\clearpage

%<PAGE-DESC:20>6 - Structural feedback e no authority inversion
%<PAGE-CONTENT:20>
\PageHashFooter{\PageMDfive{20}}
\subsection{Quando la BA espone una pressione strutturale documentale}

Durante la normalizzazione puo' emergere che un branch documentale contiene un significato reale e source-supported ma non possiede conseguenze downstream sufficientemente autonome da giustificare la propria esistenza come branch separato. Questo e' un risultato utile della BA, ma non autorizza la BA a riscrivere la documentazione.

\begin{ddtaworking}[title={Finding di review strutturale}]
Usare \PossibleStructuralRedundancy{} come segnale di review quando il significato di un branch puo' essere preservato senza perdita dentro il branch che possiede il comportamento correlato e la separazione autonoma non mostra una utilita' downstream distinguibile.
\end{ddtaworking}

Il finding deve registrare almeno:
\begin{itemize}
  \item gli elementi documentali coinvolti;
  \item il significato governato che ciascuno porta;
  \item le distinzioni analitiche che devono comunque essere preservate;
  \item le conseguenze downstream autonome cercate;
  \item perche' la separazione appare necessaria oppure potenzialmente superflua;
  \item la review documentale da riaprire.
\end{itemize}

\begin{ddtacore}[title={No authority inversion}]
La BA puo' diagnosticare la pressione e mantenere una cardinalita' analitica non artificiale. Non puo' dichiarare da sola \texttt{MERGE}, cancellare un elemento documentale o trasferirne l'ownership. La decisione strutturale torna alla Documentation Authoring review; se la documentazione cambia, la BA viene poi revalidata o ricostruita.
\end{ddtacore}
%</PAGE-CONTENT:20>

\clearpage

%<PAGE-DESC:21>6 - DermaTriage DEC02 DEC04 come pressure case
%<PAGE-CONTENT:21>
\PageHashFooter{\PageMDfive{21}}
\subsection{Pressure case DermaTriage: dominio P1--P4 e derivazione della priorita'}

La ricostruzione DermaTriage ha mostrato il caso in forma concreta. In una revisione precedente il significato relativo alla priorita' operativa era separato in due branch:

\begin{lstlisting}
DEC-MR01-02
    adozione del dominio P1-P4
    -> FR-MR01-02-01
       la priorita' appartiene a {P1,P2,P3,P4}

DEC-MR01-04
    separazione urgenza analitica / priorita' operativa
    -> FR-MR01-04-01
       regola che deriva P1-P4 da urgenza e confidence
\end{lstlisting}

La review ha chiarito che il significato ``la priorita' operativa usa il dominio P1--P4'' e' reale e deve essere preservato. Non era quindi corretto descrivere \texttt{DEC-MR01-02} come priva di semantica. Il problema era diverso: come branch autonomo non produceva un comportamento, una verifica, una modifica indipendente o una analisi distinta rispetto al branch che governa la derivazione della priorita'.

\begin{ddtaexample}[title={Esito documentale del pressure case}]
Il dominio P1--P4 viene mantenuto dentro \texttt{DEC-MR01-04}; \texttt{FR-MR01-04-01} continua a governare il mapping operativo. Il branch \texttt{DEC-MR01-02 / FR-MR01-02-01} puo' quindi essere consolidato senza perdere il significato del dominio.
\end{ddtaexample}

Questo caso mostra due regole contemporaneamente:
\begin{enumerate}
  \item una distinzione semantica puo' restare necessaria nella BA anche dopo la riduzione del numero di branch documentali;
  \item il numero di distinzioni BA non dimostra quanti branch debbano esistere nella documentazione.
\end{enumerate}
%</PAGE-CONTENT:21>

\clearpage

%<PAGE-DESC:22>6 - Downstream Utility Test e riallineamento BA
%<PAGE-CONTENT:22>
\PageHashFooter{\PageMDfive{22}}
\subsection{Ritorno alla Documentation Authoring review}

La decisione sulla struttura documentale usa il \textbf{Downstream Utility Test} definito nella Documentation Authoring Guide:

\begin{ddtacheck}[title={Downstream Utility Test - eseguito upstream}]
Se elimino questo branch come branch autonomo e trasferisco il suo significato nella Decision che possiede i comportamenti correlati, perdo una distinzione necessaria per implementazione, verifica, modifica indipendente o analisi?
\end{ddtacheck}

Se la risposta e' positiva, la separazione puo' essere giustificata. Se e' negativa, il layer documentale riesamina il branch per \texttt{MERGE} o \texttt{REWORK}. In entrambi i casi la BA non anticipa l'esito.

\begin{lstlisting}
governed / working documentation
        -> source-first BA normalization
        -> structural-feedback finding
        -> Documentation Authoring review
        -> Downstream Utility Test
        -> KEEP / MERGE / REWORK documentale
        -> BA revalidation if documentation changed
\end{lstlisting}

\begin{ddtaprinciple}[title={Distinzione da preservare}]
``Due proposition analitiche'' e ``due branch documentali'' sono domande diverse. La BA separa i significati quando serve all'analisi; la documentazione separa i branch quando serve al lavoro sul progetto. Le due cardinalita' possono coincidere, ma non sono obbligate a farlo.
\end{ddtaprinciple}

\begin{ddtaopen}[title={Confine di questa revisione}]
La R7 mantiene il catalogo in forma cumulativa: \texttt{transfer}, \texttt{produce}, \texttt{realize}, \texttt{constrain} e \texttt{condition} sono gia' ricostruiti. Il nuovo pressure case non autorizza a promuovere automaticamente un quinto operator: la rebuild testa invece se lo storico \texttt{decisionRule} possa diventare \texttt{operatorStructure.decisionRule} della proposition \texttt{produce}. Fino a regression e checkpoint, \texttt{decisionRule} resta authority corrente e la rilocazione resta candidate non ammessa.
\end{ddtaopen}
%</PAGE-CONTENT:22>

\clearpage

%<PAGE-DESC:23>7 - transfer: definizione, smallest fact e criteri d'uso
%<PAGE-CONTENT:23>
\PageHashFooter{\PageMDfive{23}}
\section{Operatore \texttt{transfer}}
\CurrentTag

\subsection{Che cosa rappresenta}
\texttt{transfer} rappresenta la \textbf{conveyance governata} di uno o piu' contenuti da una source identificabile verso una o piu' destination identificabili. Il costrutto non descrive genericamente un collegamento, un ordine di esecuzione o una vicinanza architetturale: il suo smallest semantic fact e' che \emph{un content viene trasferito da A a B}.

\begin{ddtacore}[title={Smallest semantic fact}]
Esiste un contenuto governato \texttt{X} la cui conveyance e' asserita da una source \texttt{A} verso almeno una destination \texttt{B}.
\end{ddtacore}

\subsection{Quando usarlo}
Usare \texttt{transfer} quando la documentazione governa sufficientemente tutti e tre i ruoli obbligatori: origine, destinazione e contenuto trasferito. Il verbo concreto della documentazione puo' essere ``consegnare'', ``inviare'', ``rendere disponibile a'', ``passare'', ``trasmettere'' o un'espressione equivalente; la scelta non dipende dalla parola usata ma dal significato governato.

\subsection{Quando non usarlo}
Non accettare una proposition \texttt{transfer} se almeno uno fra source, destination e content deve essere inventato per completare la frase. La presenza di due componenti nello stesso diagramma, di una chiamata API, di una sequenza di stage o di un dato disponibile in memoria non basta da sola a dimostrare conveyance.

\begin{ddtaopen}[title={Gate minimo}]
Se la source non permette di identificare \texttt{source}, almeno una \texttt{destination} e almeno un \texttt{content}, il risultato corretto e' \texttt{INSUFFICIENT EVIDENCE} / \texttt{NOT SPECIFIED}, non un transfer completato per plausibilita'.
\end{ddtaopen}
%
%</PAGE-CONTENT:23>

\clearpage

%<PAGE-DESC:24>7 - transfer: firma canonica, ruoli e cardinalita
%<PAGE-CONTENT:24>
\PageHashFooter{\PageMDfive{24}}
\subsection{Firma canonica}
La firma corrente preservata da BA2 R3 e':

\begin{lstlisting}
transfer
  behavior    -> BAReferent [0..1]
  source      -> BAReferent [1]
  destination -> BAReferent [1..*]
  content     -> BAReferent [1..*]
\end{lstlisting}

\texttt{polarity} appartiene alla \BAP{} e resta obbligatoria secondo il contratto generale. Gli scoped modifier \texttt{condition} e \texttt{temporalScope} compaiono soltanto quando la source limita realmente l'asserzione.

\begin{tabularx}{\textwidth}{L{31mm}L{21mm}Y}
\toprule
\textbf{Role} & \textbf{Card.} & \textbf{Significato} \\
\midrule
\texttt{behavior} & 0..1 & Identita' riusabile dello specifico comportamento di transfer, solo quando il segmento deve essere target di altre proposition, vincoli, realization, confronto o analisi autonoma. \\
\texttt{source} & 1 & Origine governata della conveyance. \\
\texttt{destination} & 1..* & Uno o piu' destinatari governati del content. \\
\texttt{content} & 1..* & Significato informativo/materiale la cui conveyance e' governata. \\
\bottomrule
\end{tabularx}

\subsection{Identita' opzionale del segmento}
\texttt{behavior} non viene creato per ogni scambio. Riceve identita' \BAR{} quando il trasferimento stesso deve poter essere riusato, vincolato, realizzato, protetto, confrontato fra baseline o preso come target da piu' asserzioni.

\begin{ddtacheck}[title={Behavior identity test}]
Se elimino l'identita' dello specifico segmento di transfer, perdo la possibilita' di esprimere un significato governato che riguarda proprio il passaggio e non soltanto source, destination o content? Se no, \texttt{behavior} resta assente.
\end{ddtacheck}
%
%</PAGE-CONTENT:24>

\clearpage

%<PAGE-DESC:25>7 - transfer: identity continuity, invarianti e forbidden inference
%<PAGE-CONTENT:25>
\PageHashFooter{\PageMDfive{25}}
\subsection{Continuita' di identita' del content}
Quando un significato viene prodotto, trasferito e poi usato da un consumer senza cambiare identita' semantica, la BA riusa lo stesso \BAR{}. La posizione ``output'', ``content del transfer'' e ``input'' non crea tre identita' diverse.

\begin{lstlisting}
producer result   -> X
transfer content  -> X
consumer input    -> X
\end{lstlisting}

\begin{ddtaprinciple}[title={Identity continuity rule}]
Un nuovo BAReferent e' giustificato solo quando tra i due punti compare una trasformazione che stabilisce un significato semanticamente distinto. Cambio di variabile, serializzazione, formato fisico o collocazione non bastano da soli.
\end{ddtaprinciple}

\subsection{Invarianti}
\begin{itemize}
  \item ogni transfer accettato possiede esattamente una source e almeno una destination e un content;
  \item i roleKey sono operator-scoped: \texttt{source} significa origine della conveyance, non una classificazione universale del referent;
  \item destination multiple e content multipli sono ammessi solo quando la stessa asserzione governata li comprende;
  \item l'eventuale \texttt{behavior} identifica il segmento, non sostituisce source, destination o content;
  \item polarity negativa nega l'asserzione di transfer soltanto quando questa negazione e' realmente governata; assenza di evidence non equivale a transfer negativo.
\end{itemize}

\subsection{Forbidden inference}
Dal solo \texttt{transfer} non inferire protocollo, sincronicita', cifratura, affidabilita', successo della consegna, persistenza, ownership, storage, service consumption, transport medium, retries, acknowledgement, failure semantics o topologia intermedia. Questi significati richiedono evidence e costrutti appropriati separati.
%
%</PAGE-CONTENT:25>

\clearpage

%<PAGE-DESC:26>7 - transfer: domande di identificazione e procedura di estrazione
%<PAGE-CONTENT:26>
\PageHashFooter{\PageMDfive{26}}
\subsection{Domande da porre alla documentazione}
Le domande servono a restringere l'interpretazione; non sono un classificatore automatico.

\begin{enumerate}
  \item La source governa realmente che un contenuto passi da un'origine a un destinatario?
  \item Qual e' la \texttt{source}? E' esplicitamente identificabile senza inferenza architetturale?
  \item Qual e' o quali sono le \texttt{destination}?
  \item Qual e' o quali sono i \texttt{content} trasferiti?
  \item Source, destination e content hanno gia' un \BAR{} con la stessa identita' semantica? Se si', riusarlo.
  \item Il content mantiene la stessa identita' fra producer, transfer e consumer oppure la source governa una trasformazione semanticamente distinta?
  \item Il segmento di transfer deve essere target autonomo di property, constraint, realization, security scope, change impact o altre proposition? Se si', valutare \texttt{behavior}.
  \item La proposition vale sempre oppure solo sotto una condition o un temporal scope governato?
  \item La source afferma il transfer o lo nega esplicitamente? Non dedurre polarity negativa dall'assenza di trasferimento.
  \item Rimane qualche ruolo obbligatorio non stabilito? Se si', fermarsi: il dato resta \texttt{NOT SPECIFIED} oppure diagnosis.
\end{enumerate}

\begin{ddtacheck}[title={Procedura minima}]
Leggere il smallest governed semantic fact; identificare o riusare i BAReferent necessari; verificare source/destination/content; decidere se serve behavior; compilare una sola BAProposition minima; aggiungere scoped modifier solo se governati; infine controllare le forbidden inference.
\end{ddtacheck}
%
%</PAGE-CONTENT:26>

\clearpage

%<PAGE-DESC:27>7 - transfer: esempio generico e worked example DermaTriage corrente
%<PAGE-CONTENT:27>
\PageHashFooter{\PageMDfive{27}}
\subsection{Esempio generico minimo}
\begin{ddtaexample}[title={Documentazione inventata ma coerente}]
``Il sottosistema di acquisizione MUST consegnare il \texttt{MeasurementRecord} al motore di valutazione.''
\end{ddtaexample}

Assumendo che le tre identita' abbiano superato il normale gate BA1, la proposition minima e':

\begin{lstlisting}
transfer
  source      -> MeasurementAcquisition
  destination -> EvaluationEngine
  content     -> MeasurementRecord
  polarity    -> affirmative
\end{lstlisting}

Non serve un \texttt{behavior} referent se il segmento non deve essere riusato o target di altre proposition.

\subsection{Esempio corrente - DermaTriage T06}
Nella documentazione corrente, Stage 2 produce una \texttt{ClinicalDescription} e la rende disponibile allo Stage 3; Stage 3 usa quella stessa descrizione come input del retrieval. L'estrazione corrente e':

\begin{lstlisting}
transfer
  behavior    -> ClinicalDescriptionTransfer
  source      -> ClinicalDescriptionGeneration
  destination -> HistoricalCaseRetrieval
  content     -> ClinicalDescription
  polarity    -> affirmative
\end{lstlisting}

\texttt{ClinicalDescription} e' lo stesso BAReferent gia' usato come result della generazione e come input del retrieval. \texttt{ClinicalDescriptionTransfer} riceve identita' perche' il segmento e' oggetto autonomo della review data-path / information-contract; questa scelta appartiene all'esempio, non alla definizione universale di \texttt{transfer}.

\begin{ddtaopen}[title={Boundary example}]
``Stage B usa \texttt{X}'' non basta, da solo, a costruire \texttt{transfer} se la documentazione non stabilisce da quale source \texttt{X} arrivi. L'input puo' essere noto mentre il transfer resta \texttt{INSUFFICIENT EVIDENCE}.
\end{ddtaopen}
%
%</PAGE-CONTENT:27>

\clearpage

%<PAGE-DESC:28>7 - transfer: prompt LLM per l'identificazione controllata
%<PAGE-CONTENT:28>
\PageHashFooter{\PageMDfive{28}}
\subsection{Prompt LLM di supporto all'authoring}
Il prompt e' un ausilio operativo: non cambia la semantica BA e non rende autorevole una inferenza non supportata. Deve ricevere l'estratto documentale governato e, quando disponibile, il registro dei BAReferent gia' esistenti.

\begin{lstlisting}
TASK: identify governed BA transfer propositions.
Use ONLY the supplied documentation and existing BA registry.

A transfer exists only if the source supports:
  source [1], destination [1..*], content [1..*].

Rules:
1. Identify the smallest governed conveyance fact.
2. Reuse existing BAReferents when semantic identity is unchanged.
3. Propose a new BAReferent only when BA1 independent identity is needed.
4. Add behavior only if the transfer segment itself must be independently
   targeted, constrained, realized, compared, traced or analyzed.
5. Preserve the same content BAReferent across producer/transfer/consumer
   when the governed meaning is unchanged.
6. Use polarity=negative only for an explicitly negated transfer assertion.
7. Add condition/temporalScope only when explicitly governed.
8. If source, destination or content is missing, return INSUFFICIENT_EVIDENCE.
9. Do not infer protocol, network topology, persistence, success, ownership,
   encryption, service consumption, transport medium or failure semantics.
10. Do not invent operators, roles, states or metadata.

Return EXACTLY the authoring wrapper and BA output shape shown next.
\end{lstlisting}

\begin{ddtaprinciple}[title={Human review remains mandatory}]
L'LLM puo' proporre referent e proposition; l'accettazione richiede comunque verifica contro source, identity criterion, canonical naming e provenance. Una risposta sintatticamente completa non sana un gap documentale.
\end{ddtaprinciple}
%
%</PAGE-CONTENT:28>

\clearpage

%<PAGE-DESC:29>7 - transfer: output BA atteso e discriminazioni differite
%<PAGE-CONTENT:29>
\PageHashFooter{\PageMDfive{29}}
\subsection{Output esatto atteso dall'estrazione assistita}
Il wrapper seguente e' un formato di authoring/review, non un nuovo costrutto BA. La parte di output BA usa soltanto la shape ammessa.

\begin{lstlisting}
TRANSFER_ASSESSMENT
  decision = ACCEPT_TRANSFER | NOT_A_TRANSFER | INSUFFICIENT_EVIDENCE
  evidence = <source locator / concise supported wording>
  notSpecified = [<only unresolved material items>]

BA_OUTPUT
  BAReferent reuse/candidates:
    <canonical names only as needed>

  BAProposition <candidate-id>
    semanticOperatorKey = transfer
    behavior    -> <BAReferent>      # omit line if not justified
    source      -> <BAReferent>
    destination -> <BAReferent>      # repeat for 1..*
    content     -> <BAReferent>      # repeat for 1..*
    polarity    = affirmative | negative
    scopedModifier = <condition|temporalScope>  # omit if absent
\end{lstlisting}

Se \texttt{decision != ACCEPT\_TRANSFER}, \texttt{BA\_OUTPUT} non contiene una proposition transfer accettabile. I nomi proposti restano candidate finche' non superano canonical naming, provenance e review.

\subsection{Discriminazioni con altri costrutti}
-
%
%</PAGE-CONTENT:29>

\clearpage

%<PAGE-DESC:30>7 - transfer: rappresentazione Mermaid
%<PAGE-CONTENT:30>
\PageHashFooter{\PageMDfive{30}}
\subsection{Rappresentazione Mermaid di transfer}

Per \texttt{transfer}, la projection rende direttamente leggibili i role dell'operatore. \texttt{source} e \texttt{destination} sono gli endpoint; \texttt{content} occupa la posizione centrale ed e' mostrato con una forma documento. L'intero percorso direzionale \texttt{source -> content -> destination} rappresenta \textbf{una sola} \BAP{} \texttt{transfer}: le due frecce sono una decomposizione grafica del percorso, non due proposition.

Quando il role opzionale \texttt{behavior} e' presente, viene visualizzato come annotazione distinta del transfer: nodo giallo collegato al content con una linea tratteggiata e non direzionale. Questa linea e' una convenzione di projection del role \texttt{behavior}; non rappresenta una seconda relazione BA.

\begin{tabularx}{\textwidth}{L{42mm}Y}
\toprule
\textbf{Parte di transfer} & \textbf{Rappresentazione Mermaid} \\
\midrule
\texttt{source} & Nodo endpoint all'origine del percorso. \\
\texttt{content} & Nodo centrale con forma documento e sfondo bianco. \\
\texttt{destination} & Nodo endpoint raggiunto dal percorso direzionale. \\
\texttt{behavior} & Nodo opzionale con sfondo giallo; linea tratteggiata non direzionale verso il content. \\
Proposition id / polarity & Metadati della stessa \BAP{}; non diventano un nodo del data path. \\
\bottomrule
\end{tabularx}

\begin{lstlisting}
flowchart LR
  S["ClinicalDescriptionGeneration<br/>source"]
  C@{ shape: doc, label: "ClinicalDescription<br/>content" }
  D["HistoricalCaseRetrieval<br/>destination"]
  B["ClinicalDescriptionTransfer<br/>behavior"]

  S --> C --> D
  B --- C

  classDef endpoint fill:#EEF4F9,stroke:#275D86,color:#17324D
  classDef content fill:#FFFFFF,stroke:#275D86,color:#17324D
  classDef behavior fill:#FFF4E3,stroke:#A96A16,color:#17324D
  class S,D endpoint
  class C content
  class B behavior
  linkStyle 2 stroke:#A96A16,stroke-width:1.5px,stroke-dasharray:5 5
\end{lstlisting}

Il codice usa il caso T06 come esempio concreto. Se \texttt{behavior} fosse assente, il nodo giallo e la linea tratteggiata verrebbero omessi; il percorso \texttt{source -> content -> destination} resterebbe invariato.
%</PAGE-CONTENT:30>

\clearpage

%<PAGE-DESC:31>7 - transfer: esempio grafico Mermaid
%<PAGE-CONTENT:31>
\PageHashFooter{\PageMDfive{31}}
\subsection{Esempio grafico corrispondente}

Il diagramma seguente applica la stessa convenzione al caso T06: il documento bianco e' il \texttt{content}; il percorso con frecce rende la direzione del \texttt{transfer}; il nodo giallo identifica il \texttt{behavior} opzionale senza trasformarlo in uno stadio del flusso.

\vspace{5mm}
\begin{center}
\begin{tikzpicture}[
  endpoint/.style={draw=DDTABlue,fill=DDTALightBlue,rounded corners=3pt,thick,align=center,minimum height=18mm,text width=37mm,inner sep=6pt,font=\sffamily\fontsize{7}{8.2}\selectfont},
  contentdoc/.style={draw=DDTABlue,fill=white,thick,align=center,minimum height=22mm,text width=32mm,inner sep=7pt,font=\sffamily\fontsize{7}{8.2}\selectfont,
    path picture={\draw[DDTABlue,thick] ([xshift=-6mm]path picture bounding box.north east) -- ([xshift=-6mm,yshift=-6mm]path picture bounding box.north east) -- ([yshift=-6mm]path picture bounding box.north east); \draw[DDTABlue,thick] ([xshift=-6mm]path picture bounding box.north east) -- ([yshift=-6mm]path picture bounding box.north east);}},
  behavior/.style={draw=DDTAAmber,fill=DDTALightAmber,rounded corners=3pt,thick,align=center,minimum height=15mm,text width=40mm,inner sep=6pt,font=\sffamily\fontsize{7}{8.2}\selectfont},
  flow/.style={draw=DDTABlue,very thick,-{Stealth[length=3mm,width=2mm]}},
  qualifier/.style={draw=DDTAAmber,thick,dashed}
]
\node[contentdoc] (c) {\textbf{ClinicalDescription}\\{\scriptsize\texttt{content}}};
\node[endpoint,left=15mm of c] (s) {\textbf{ClinicalDescriptionGeneration}\\{\scriptsize\texttt{source}}};
\node[endpoint,right=15mm of c] (d) {\textbf{HistoricalCaseRetrieval}\\{\scriptsize\texttt{destination}}};
\node[behavior,above=17mm of c] (b) {\textbf{ClinicalDescriptionTransfer}\\{\scriptsize\texttt{behavior}}};

\draw[flow] (s) -- (c);
\draw[flow] (c) -- (d);
\draw[qualifier] (b) -- (c);
\end{tikzpicture}
\end{center}

\vspace{3mm}
\begin{center}
\sffamily\small\textbf{BAP-T06-01}\quad \texttt{transfer}\quad \texttt{polarity = affirmative}
\end{center}
\vspace{3mm}

\begin{ddtacheck}[title={Lettura del grafico}]
Il percorso completo da \texttt{source} a \texttt{destination}, passando per il documento \texttt{content}, e' la projection di una sola proposition \texttt{transfer}. La linea tratteggiata verso il nodo giallo esprime soltanto la presenza del role opzionale \texttt{behavior}; non e' una BAProposition aggiuntiva e non introduce una seconda direzione di flusso.
\end{ddtacheck}
%</PAGE-CONTENT:31>


\clearpage

%<PAGE-DESC:32>8 - produce: definizione, smallest fact e criteri d'uso
%<PAGE-CONTENT:32>
\PageHashFooter{\PageMDfive{32}}
\section{Operatore \texttt{produce}}
\CurrentTag

\subsection{Che cosa rappresenta}
\texttt{produce} rappresenta l'asserzione governata che un actor o una capability \textbf{rende disponibile uno o piu' risultati/output}, eventualmente usando input esplicitamente governati. Il costrutto descrive il fatto produttivo minimo; non descrive automaticamente come il risultato viene calcolato, dove viene memorizzato o a chi viene successivamente trasferito.

\begin{ddtacore}[title={Smallest semantic fact}]
Esiste un actor/capability \texttt{A} che rende disponibile almeno un risultato governato \texttt{R}. Eventuali input \texttt{I} appartengono alla stessa proposition solo quando la documentazione governa esplicitamente che vengano usati per quella produzione.
\end{ddtacore}

\subsection{Quando usarlo}
Usare \texttt{produce} quando la documentazione governa che un actor/capability produce, genera, calcola, restituisce, emette o rende disponibile un risultato identificabile. Il verbo concreto non e' decisivo: deve essere governato il significato ``A rende disponibile R''. Gli input sono opzionali e vengono inclusi soltanto quando la source li collega alla produzione del risultato.

\subsection{Quando non usarlo}
Non accettare una proposition \texttt{produce} quando il risultato deve essere inventato, quando e' nota soltanto l'esistenza di un dato senza un actor che lo renda disponibile, oppure quando un input viene aggiunto solo perche' compare altrove nell'architettura. La presenza di una sequenza di stage, di un componente o di una dipendenza temporale non dimostra da sola una produzione.

\begin{ddtaopen}[title={Gate minimo}]
Se la source non permette di identificare esattamente un \texttt{actor} e almeno un \texttt{result}, il risultato corretto e' \texttt{INSUFFICIENT EVIDENCE} / \texttt{NOT SPECIFIED}. L'assenza di input espliciti, invece, non invalida \texttt{produce}: \texttt{input} ha cardinalita' \texttt{0..*}.
\end{ddtaopen}
%
%</PAGE-CONTENT:32>

\clearpage

%<PAGE-DESC:33>8 - produce: firma canonica, ruoli e cardinalita
%<PAGE-CONTENT:33>
\PageHashFooter{\PageMDfive{33}}
\subsection{Firma canonica}
La firma completa e':

\begin{lstlisting}
produce
  actor    -> BAReferent [1]
  input    -> BAReferent [0..*]
  result   -> BAReferent [1..*]
  polarity = affirmative | negative [1]
\end{lstlisting}

Gli scoped modifier \texttt{condition} e \texttt{temporalScope} compaiono soltanto quando la source limita realmente l'asserzione.

\begin{tabularx}{\textwidth}{L{31mm}L{21mm}Y}
\toprule
\textbf{Role} & \textbf{Card.} & \textbf{Significato} \\
\midrule
\texttt{actor} & 1 & Actor/capability a cui la documentazione attribuisce il fatto di rendere disponibile il risultato. \\
\texttt{input} & 0..* & Significati che la documentazione dichiara usati dalla stessa produzione; non vengono ricostruiti da dipendenze plausibili. \\
\texttt{result} & 1..* & Uno o piu' risultati/output governati resi disponibili dall'actor. \\
\bottomrule
\end{tabularx}

%
%</PAGE-CONTENT:33>

\clearpage

%<PAGE-DESC:34>8 - produce: identity continuity, invarianti e forbidden inference
%<PAGE-CONTENT:34>
\PageHashFooter{\PageMDfive{34}}
\subsection{Continuita' di identita' di input e result}
Un \texttt{result} non riceve una nuova identita' soltanto perche' viene successivamente trasferito o usato come \texttt{input} da un altro comportamento. Quando il significato governato resta lo stesso, la BA riusa lo stesso \BAR{}.

\begin{lstlisting}
producer result   -> X
transfer content  -> X
consumer input    -> X
\end{lstlisting}

Allo stesso modo, un \texttt{input} di \texttt{produce} e' il medesimo significato gia' identificato altrove quando non esiste una trasformazione semanticamente distinta.

\subsection{Invarianti}
\begin{itemize}
  \item ogni \texttt{produce} accettato possiede esattamente un \texttt{actor} e almeno un \texttt{result};
  \item \texttt{input} e' opzionale: la proposition resta valida anche quando la source non governa input espliciti;
  \item input multipli o result multipli appartengono alla stessa proposition solo quando la stessa asserzione governata li lega allo stesso fatto produttivo;
  \item i roleKey sono operator-scoped: essere \texttt{result} in una proposition non assegna al referent un tipo universale ``output'';
  \item polarity negativa richiede una negazione governata della produzione; assenza di evidence non equivale a \texttt{produce} negativo.
\end{itemize}

\subsection{Forbidden inference}
Dal solo \texttt{produce} non inferire algoritmo, modello, trasformazioni intermedie, destination, transfer, protocollo, serializzazione, persistence, storage, ownership, creazione ex novo dell'item, successo runtime, retry, failure semantics o completezza dello schema del risultato. Non aggiungere inoltre un \texttt{input} soltanto perche' un dato e' disponibile nel sistema o compare in uno stadio precedente.
%
%</PAGE-CONTENT:34>

\clearpage

%<PAGE-DESC:35>8 - produce: domande di identificazione e procedura di estrazione
%<PAGE-CONTENT:35>
\PageHashFooter{\PageMDfive{35}}
\subsection{Domande da porre alla documentazione}
Le domande servono a restringere l'interpretazione; non sono un classificatore automatico.

\begin{enumerate}
  \item La source governa realmente che qualcosa renda disponibile un risultato/output?
  \item Qual e' esattamente l'\texttt{actor}? E' attribuito dalla documentazione senza inferenza architetturale?
  \item Qual e' o quali sono i \texttt{result} prodotti?
  \item La source governa uno o piu' \texttt{input} usati per quella produzione? Se non li governa, lasciarli assenti.
  \item Actor, input e result hanno gia' un \BAR{} con la stessa identita' semantica? Se si', riusarlo.
  \item Piu' input o piu' result appartengono davvero allo stesso fatto produttivo oppure la documentazione governa produzioni distinte?
  \item Un result mantiene la stessa identita' quando viene trasferito o riusato downstream? Se si', non creare un nuovo referent per il cambio di posizione.
  \item La proposition vale sempre oppure solo sotto una condition o un temporal scope governato?
  \item La source afferma la produzione o la nega esplicitamente? Non dedurre polarity negativa dall'assenza del risultato.
  \item Rimane non stabilito l'actor o ogni result? Se si', fermarsi: il dato resta \texttt{NOT SPECIFIED} oppure diagnosis.
\end{enumerate}

\begin{ddtacheck}[title={Procedura minima}]
Leggere il smallest governed semantic fact; identificare o riusare l'actor e almeno un result; aggiungere soltanto gli input esplicitamente legati alla produzione; compilare una sola BAProposition minima; aggiungere scoped modifier solo se governati; infine controllare le forbidden inference e la continuita' di identita' downstream.
\end{ddtacheck}
%
%</PAGE-CONTENT:35>

\clearpage

%<PAGE-DESC:36>8 - produce: esempio generico e worked example DermaTriage corrente
%<PAGE-CONTENT:36>
\PageHashFooter{\PageMDfive{36}}
\subsection{Esempio generico minimo}
\begin{ddtaexample}[title={Documentazione inventata ma coerente}]
``La capability \texttt{RiskAssessmentComputation} MUST usare il \texttt{MeasurementRecord} per rendere disponibili un \texttt{RiskAssessment} e una \texttt{AssessmentConfidence}.''
\end{ddtaexample}

Assumendo che le identita' abbiano superato il normale gate BA1, la proposition e':

\begin{lstlisting}
produce
  actor  -> RiskAssessmentComputation
  input  -> MeasurementRecord
  result -> RiskAssessment
  result -> AssessmentConfidence
  polarity -> affirmative
\end{lstlisting}

I due result restano nella stessa proposition perche' la frase governa un unico fatto produttivo che li rende disponibili insieme.

\subsection{Esempio corrente - DermaTriage Stage 3}
Nella documentazione corrente, lo Stage 3 usa la \texttt{ClinicalDescription} prodotta dallo Stage 2 per il retrieval e rende disponibile il contesto dei casi storici recuperati. La proposition accettata corrente e':

\begin{lstlisting}
BAP-FR03-03-01
semanticOperatorKey = produce
  actor  -> HistoricalCaseRetrieval
  input  -> ClinicalDescription
  result -> RetrievedHistoricalCaseContext
polarity = affirmative
\end{lstlisting}

\texttt{ClinicalDescription} e' lo stesso \BAR{} gia' usato come \texttt{content} del transfer T06; non viene duplicato come ``Stage3Input''. Il dettaglio ChromaDB / embedding appartiene ad altre proposition gia' separate e non viene incorporato nella shape di \texttt{produce}.

\begin{ddtaopen}[title={Boundary example}]
``Il componente B dispone di \texttt{X}'' non basta, da solo, a costruire \texttt{produce}: puo' descrivere disponibilita' o stato senza governare chi produce \texttt{X}. Se actor o result non sono stabiliti, l'esito resta \texttt{INSUFFICIENT EVIDENCE}.
\end{ddtaopen}
%
%</PAGE-CONTENT:36>

\clearpage

%<PAGE-DESC:37>8 - produce: prompt LLM per l'identificazione controllata
%<PAGE-CONTENT:37>
\PageHashFooter{\PageMDfive{37}}
\subsection{Prompt LLM di supporto all'authoring}
Il prompt e' un ausilio operativo: non cambia la semantica BA e non rende autorevole una inferenza non supportata. Deve ricevere l'estratto documentale governato e, quando disponibile, il registro dei BAReferent gia' esistenti.

\begin{lstlisting}
TASK: identify governed BA produce propositions.
Use ONLY the supplied documentation and existing BA registry.

A produce proposition exists only if the source supports:
  actor [1], result [1..*].
Input is optional: input [0..*].

Rules:
1. Identify the smallest governed production fact: actor makes result available.
2. Reuse existing BAReferents when semantic identity is unchanged.
3. Propose a new BAReferent only when BA1 independent identity is needed.
4. Add input only when the source explicitly governs its use in this production.
5. Keep multiple inputs/results in one proposition only when the same governed
   production fact binds them to the same actor.
6. Preserve result identity when it later appears as transfer content or input.
7. Use polarity=negative only for an explicitly negated production assertion.
8. Add condition/temporalScope only when explicitly governed.
9. If actor or every result is missing, return INSUFFICIENT_EVIDENCE.
10. Do not infer algorithm, implementation, destination, transfer, persistence,
    storage, creation, ownership, serialization or failure semantics.
11. Do not invent operators, roles, states or metadata.

Return EXACTLY the authoring wrapper and BA output shape shown next.
\end{lstlisting}

\begin{ddtaprinciple}[title={Human review remains mandatory}]
L'LLM puo' proporre referent e proposition; l'accettazione richiede comunque verifica contro source, identity criterion, canonical naming e provenance. Una risposta sintatticamente completa non sana un gap documentale e non rende implicito un input non governato.
\end{ddtaprinciple}
%
%</PAGE-CONTENT:37>

\clearpage

%<PAGE-DESC:38>8 - produce: output BA atteso e discriminazioni differite
%<PAGE-CONTENT:38>
\PageHashFooter{\PageMDfive{38}}
\subsection{Output esatto atteso dall'estrazione assistita}
Il wrapper seguente e' un formato di authoring/review, non un nuovo costrutto BA. La parte di output BA usa soltanto la shape ammessa.

\begin{lstlisting}
PRODUCE_ASSESSMENT
  decision = ACCEPT_PRODUCE | NOT_A_PRODUCE | INSUFFICIENT_EVIDENCE
  evidence = <source locator / concise supported wording>
  notSpecified = [<only unresolved material items>]

BA_OUTPUT
  BAReferent reuse/candidates:
    <canonical names only as needed>

  BAProposition <candidate-id>
    semanticOperatorKey = produce
    actor  -> <BAReferent>
    input  -> <BAReferent>          # repeat for 0..*; omit if absent
    result -> <BAReferent>          # repeat for 1..*
    polarity = affirmative | negative
    scopedModifier = <condition|temporalScope>  # omit if absent
\end{lstlisting}

Se \texttt{decision != ACCEPT\_PRODUCE}, \texttt{BA\_OUTPUT} non contiene una proposition \texttt{produce} accettabile. I nomi proposti restano candidate finche' non superano canonical naming, provenance e review. L'assenza di \texttt{input} non e' un errore quando la source governa actor e result ma non esplicita gli input.

\subsection{Discriminazioni con altri costrutti}
-
%
%</PAGE-CONTENT:38>

\clearpage

%<PAGE-DESC:39>8 - produce: rappresentazione Mermaid
%<PAGE-CONTENT:39>
\PageHashFooter{\PageMDfive{39}}
\subsection{Rappresentazione Mermaid di produce}

Per \texttt{produce}, la projection mette l'\texttt{actor} al centro del fatto produttivo. Gli eventuali \texttt{input} entrano nell'actor; i \texttt{result} escono dall'actor. L'intero insieme di frecce associato allo stesso proposition id rappresenta \textbf{una sola} \BAP{} \texttt{produce}: le singole frecce visualizzano le participation della proposition e non creano proposition aggiuntive.

Nel worked example DermaTriage, \texttt{ClinicalDescription} e \texttt{RetrievedHistoricalCaseContext} sono mostrati con lo stesso simbolo documento bianco gia' usato per il contenuto informativo nella projection di \texttt{transfer}. Il role resta pero' quello della proposition corrente: rispettivamente \texttt{input} e \texttt{result}. La forma grafica non rinomina i role BA.

\begin{tabularx}{\textwidth}{L{42mm}Y}
\toprule
\textbf{Parte di produce} & \textbf{Rappresentazione Mermaid} \\
\midrule
\texttt{input} & Nodo opzionale a sinistra; nel worked example informativo usa il simbolo documento bianco. Freccia direzionale verso l'actor. Ripetibile per \texttt{0..*}. \\
\texttt{actor} & Nodo centrale che identifica l'actor/capability della production. \\
\texttt{result} & Nodo a destra; nel worked example informativo usa il simbolo documento bianco. Freccia in uscita dall'actor. Ripetibile per \texttt{1..*}. \\
Proposition id / polarity & Metadati dell'intero pattern; non diventano nodi del flusso. \\
\bottomrule
\end{tabularx}

\begin{lstlisting}
flowchart LR
  I@{ shape: doc, label: "ClinicalDescription<br/>input" }
  A["HistoricalCaseRetrieval<br/>actor"]
  R@{ shape: doc, label: "RetrievedHistoricalCaseContext<br/>result" }

  I --> A --> R

  classDef actor fill:#EEF4F9,stroke:#275D86,color:#17324D
  classDef content fill:#FFFFFF,stroke:#275D86,color:#17324D
  class A actor
  class I,R content
\end{lstlisting}

Se \texttt{input} fosse assente, il nodo a sinistra e la relativa freccia verrebbero omessi. Con piu' input o result, i nodi vengono ripetuti e convergono/divergono dall'unico actor mantenendo lo stesso proposition id.
%
%</PAGE-CONTENT:39>

\clearpage

%<PAGE-DESC:40>8 - produce: esempio grafico Mermaid
%<PAGE-CONTENT:40>
\PageHashFooter{\PageMDfive{40}}
\subsection{Esempio grafico corrispondente}

Il diagramma seguente applica la convenzione alla proposition DermaTriage \texttt{BAP-FR03-03-01}. Il nodo centrale e' l'\texttt{actor}; l'input entra nell'actor e il result ne esce. I due referent informativi mantengono il simbolo documento bianco usato nella projection precedente. Il grafico non aggiunge destination, storage o dettagli di implementazione che \texttt{produce} non governa.

\vspace{7mm}
\begin{center}
\begin{tikzpicture}[
  actor/.style={draw=DDTABlue,fill=DDTALightBlue,rounded corners=3pt,thick,align=center,minimum height=20mm,text width=43mm,inner sep=7pt,font=\sffamily\fontsize{7}{8.2}\selectfont},
  contentdoc/.style={draw=DDTABlue,fill=white,thick,align=center,minimum height=22mm,text width=40mm,inner sep=7pt,font=\sffamily\fontsize{7}{8.2}\selectfont,
    path picture={\draw[DDTABlue,thick] ([xshift=-6mm]path picture bounding box.north east) -- ([xshift=-6mm,yshift=-6mm]path picture bounding box.north east) -- ([yshift=-6mm]path picture bounding box.north east); \draw[DDTABlue,thick] ([xshift=-6mm]path picture bounding box.north east) -- ([yshift=-6mm]path picture bounding box.north east);}},
  flow/.style={draw=DDTABlue,very thick,-{Stealth[length=3mm,width=2mm]}}
]
\node[actor] (a) {\textbf{HistoricalCaseRetrieval}\\{\scriptsize\texttt{actor}}};
\node[contentdoc,left=16mm of a] (i) {\textbf{ClinicalDescription}\\{\scriptsize\texttt{input}}};
\node[contentdoc,right=16mm of a] (r) {\textbf{RetrievedHistoricalCaseContext}\\{\scriptsize\texttt{result}}};

\draw[flow] (i) -- (a);
\draw[flow] (a) -- (r);
\end{tikzpicture}
\end{center}

\vspace{5mm}
\begin{center}
\sffamily\small\textbf{BAP-FR03-03-01}\quad \texttt{produce}\quad \texttt{polarity = affirmative}
\end{center}
\vspace{4mm}

\begin{ddtacheck}[title={Lettura del grafico}]
L'intero pattern \texttt{input -> actor -> result} e' la projection di una sola proposition \texttt{produce}. Se la stessa proposition avesse piu' input o piu' result, comparirebbero piu' rami verso o dall'actor senza cambiare il numero di BAProposition rappresentate.
\end{ddtacheck}
%</PAGE-CONTENT:40>
\clearpage


%<PAGE-DESC:41>9 - realize: definizione, smallest fact e criteri d'uso
%<PAGE-CONTENT:41>
\PageHashFooter{\PageMDfive{41}}
\section{Operatore \texttt{realize}}
\CurrentTag

\subsection{Che cosa rappresenta}
\texttt{realize} rappresenta l'asserzione governata che uno o piu' significati concreti \textbf{realizzano o materializzano un significato piu' astratto} nella baseline considerata. Il costrutto preserva la distinzione tra \emph{che cosa} il progetto richiede o governa a livello astratto e \emph{quale elemento concreto} ne costituisce la realizzazione corrente.

\begin{ddtacore}[title={Smallest semantic fact}]
Esiste un significato astratto \texttt{A} e almeno una realizzazione concreta \texttt{R} tale che la documentazione governa che \texttt{R} realizzi/materializzi \texttt{A}. I due significati restano distinti BAReferent: la relation non li fonde e non trasforma automaticamente l'abstract in una classe software.
\end{ddtacore}

\subsection{Quando usarlo}
Usare \texttt{realize} quando la documentazione distingue un significato astratto --- per esempio capability, contract, behavior o strategia richiesta --- da uno o piu' elementi concreti scelti per materializzarlo nella baseline. La relazione deve essere governata: la semplice presenza di una tecnologia dentro un componente non basta.

\subsection{Quando non usarlo}
Non usare \texttt{realize} per significare che un componente usa una libreria, legge dati da un database, dipende da un servizio, produce un output o trasferisce informazione. Questi fatti possono richiedere altre proposition e non diventano realization solo perche' il concreto compare nell'implementazione.

\begin{ddtaopen}[title={Gate minimo}]
Se la source identifica soltanto l'elemento concreto ma non il significato astratto che esso realizza, non inventare retroattivamente l'abstract per completare la BA. Registrare \texttt{INSUFFICIENT EVIDENCE} / \texttt{NOT SPECIFIED} e, quando materialmente utile, riportare il gap alla review documentale.
\end{ddtaopen}
%</PAGE-CONTENT:41>

\clearpage

%<PAGE-DESC:42>9 - realize: firma canonica, ruoli e cardinalita
%<PAGE-CONTENT:42>
\PageHashFooter{\PageMDfive{42}}
\subsection{Firma canonica}
La firma completa e':

\begin{lstlisting}
realize
  abstract     -> BAReferent [1]
  realization  -> BAReferent [1..*]
  polarity     = affirmative | negative [1]
\end{lstlisting}

Gli scoped modifier \texttt{condition} e \texttt{temporalScope} compaiono soltanto quando la source limita realmente l'asserzione.

\begin{tabularx}{\textwidth}{L{31mm}L{21mm}Y}
\toprule
\textbf{Role} & \textbf{Card.} & \textbf{Significato} \\
\midrule
\texttt{abstract} & 1 & Significato astratto che la baseline governa indipendentemente dalla realizzazione concreta scelta. \\
\texttt{realization} & 1..* & Uno o piu' BAReferent concreti che la documentazione identifica come materializzazione dell'abstract. \\
\bottomrule
\end{tabularx}

\begin{ddtaprinciple}[title={Cardinalita' multipla non significa alternativa}]
La cardinalita' \texttt{realization [1..*]} permette a una proposition di riferirsi a piu' realization governate. Da questa cardinalita' non si deduce che esse siano alternative, intercambiabili, ridondanti o usate congiuntamente: tali distinzioni richiedono evidence separata.
\end{ddtaprinciple}
%</PAGE-CONTENT:42>

\clearpage

%<PAGE-DESC:43>9 - realize: identity continuity, invarianti e forbidden inference
%<PAGE-CONTENT:43>
\PageHashFooter{\PageMDfive{43}}
\subsection{Continuita' di identita' tra abstract e realization}
\texttt{abstract} e \texttt{realization} sono significati distinti e mantengono identita' indipendente. La relation serve proprio a preservare contemporaneamente il significato stabile e la scelta concreta della baseline.

\begin{lstlisting}
abstract capability  -> ClinicalDescriptionGeneration
realization concrete -> Qwen2-VL-7B-Instruct
\end{lstlisting}

La realization non sostituisce il BAReferent astratto. Se altri documenti qualificano, vincolano o riusano \texttt{ClinicalDescriptionGeneration}, tali proposition continuano a riferirsi all'abstract e non vengono automaticamente trasferite alla tecnologia concreta.

\subsection{Invarianti}
\begin{itemize}
  \item ogni \texttt{realize} accettato possiede esattamente un \texttt{abstract} e almeno una \texttt{realization};
  \item abstract e realization restano BAReferent distinti e riusabili;
  \item piu' realization appartengono alla stessa proposition soltanto quando la stessa asserzione governata le lega allo stesso abstract;
  \item la relation deve esprimere materializzazione del significato astratto, non mera presenza tecnica o collaborazione;
  \item polarity negativa richiede una negazione governata della realization; assenza di evidence non equivale a \texttt{realize} negativo.
\end{itemize}

\subsection{Forbidden inference}
Dal solo \texttt{realize} non inferire subclassing, ereditarieta' software, equivalenza semantica completa, sostituibilita' drop-in, compatibilita' di interfaccia, deployment topology, ownership, dipendenza, consumo di servizio, flusso dati, produzione, transfer, algoritmo interno, modello di sicurezza o completezza delle realization possibili. La projection puo' ricordare una gerarchia abstract/concrete, ma non importa automaticamente la semantica UML delle classi.
%</PAGE-CONTENT:43>

\clearpage

%<PAGE-DESC:44>9 - realize: domande di identificazione e procedura di estrazione
%<PAGE-CONTENT:44>
\PageHashFooter{\PageMDfive{44}}
\subsection{Domande da porre alla documentazione}
Le domande cercano la distinzione abstract/concrete prima della sintassi BA.

\begin{enumerate}
  \item Qual e' il significato astratto che la documentazione governa indipendentemente dalla tecnologia o strategia corrente?
  \item Quale elemento concreto la source dichiara che realizza, implementa o materializza quel significato?
  \item Se rimuovo il nome della tecnologia concreta, l'abstract resta un significato comprensibile e governato? Se no, verificare che non stiamo inventando un livello astratto.
  \item Il concreto realizza realmente l'abstract oppure lo usa, lo invoca, ne dipende, gli fornisce input o ne riceve output?
  \item Abstract e realization hanno gia' BAReferent con la stessa identita' semantica? Se si', riusarli.
  \item Se compaiono piu' realization, la documentazione le lega davvero allo stesso abstract nella stessa asserzione? Non inferire alternative o composizione.
  \item La source governa condition o temporal scope della realization?
  \item La source afferma o nega esplicitamente la realization? Non dedurre polarity negativa dall'assenza di una tecnologia.
  \item La tecnologia e' nominata ma la funzione astratta resta implicita? Fermarsi e registrare un possibile gap documentale invece di completarlo per plausibilita'.
\end{enumerate}

\begin{ddtacheck}[title={Procedura minima}]
Isolare prima il significato astratto; verificare poi che la source identifichi esplicitamente almeno una realizzazione concreta; riusare le identita' BA esistenti; compilare una sola proposition minima; infine riesaminare che nessuna relazione di uso, dipendenza, produzione o trasferimento sia stata compressa impropriamente dentro \texttt{realize}.
\end{ddtacheck}
%</PAGE-CONTENT:44>

\clearpage

%<PAGE-DESC:45>9 - realize: esempio generico e worked example DermaTriage corrente
%<PAGE-CONTENT:45>
\PageHashFooter{\PageMDfive{45}}
\subsection{Esempio generico minimo}
\begin{ddtaexample}[title={Documentazione inventata ma coerente}]
``La capability \texttt{DocumentClassification} MUST essere realizzata nella baseline corrente dalla strategia \texttt{TransformerClassificationStrategy}.''
\end{ddtaexample}

Assumendo che le identita' abbiano superato il normale gate BA1, la proposition e':

\begin{lstlisting}
realize
  abstract     -> DocumentClassification
  realization  -> TransformerClassificationStrategy
  polarity     -> affirmative
\end{lstlisting}

Il fatto preservato e' la scelta concreta rispetto a un significato astratto gia' distinguibile; non vengono inferiti modello specifico, formato degli input o contratto degli output.

\subsection{Esempio corrente - DermaTriage Stage 2}
La documentazione corrente distingue la capability \texttt{ClinicalDescriptionGeneration} dalla realizzazione concreta usata nello Stage 2. L'analisi corrente rappresenta quindi:

\begin{lstlisting}
BAP-FR03-02-02
semanticOperatorKey = realize
  abstract     -> ClinicalDescriptionGeneration
  realization  -> Qwen2-VL-7B-Instruct
polarity = affirmative
\end{lstlisting}

Questa proposition non sostituisce la proposition \texttt{produce} che governa la produzione di \texttt{ClinicalDescription}. \texttt{realize} risponde a ``quale elemento concreto materializza la capability?'', mentre \texttt{produce} preserva il fatto che la capability rende disponibile il risultato.

\begin{ddtaopen}[title={Boundary example}]
``Il componente \texttt{A} usa \texttt{TechnologyB}'' non basta a costruire \texttt{realize}. Senza un abstract governato e una relazione di materializzazione esplicita, l'uso della tecnologia resta un fatto diverso oppure un gap da riesaminare nella documentazione.
\end{ddtaopen}
%</PAGE-CONTENT:45>

\clearpage

%<PAGE-DESC:46>9 - realize: prompt LLM per l'identificazione controllata
%<PAGE-CONTENT:46>
\PageHashFooter{\PageMDfive{46}}
\subsection{Prompt LLM di supporto all'authoring}
Il prompt non deve creare un abstract soltanto per rendere sostituibile o piu' elegante una tecnologia gia' nominata. Un gap documentale resta un gap.

\begin{lstlisting}
TASK: identify governed BA realize propositions.
Use ONLY the supplied documentation and existing BA registry.

A realize proposition exists only if the source supports:
  abstract [1], realization [1..*].

Rules:
1. Identify the smallest governed abstract/concrete materialization fact.
2. Reuse existing BAReferents when semantic identity is unchanged.
3. Propose a new BAReferent only when BA1 independent identity is needed.
4. Accept realization only when the source governs that the concrete referent
   realizes/materializes the abstract referent.
5. Do not convert mere use, dependency, invocation, production or data flow
   into realization.
6. Keep multiple realizations in one proposition only when the same governed
   assertion binds them to the same abstract.
7. Do not infer alternatives, replaceability, composition or completeness.
8. Use polarity=negative only for an explicitly negated realization assertion.
9. Add condition/temporalScope only when explicitly governed.
10. If abstract or every realization is missing, return INSUFFICIENT_EVIDENCE.
11. If a concrete technology is named but its abstract function is not governed,
    report the documentation gap; do not invent the abstraction.
12. Do not invent operators, roles, states or metadata.

Return EXACTLY the authoring wrapper and BA output shape shown next.
\end{lstlisting}

\begin{ddtaprinciple}[title={Human review remains mandatory}]
La distinzione fra ``realizza'' e ``usa'' e' semanticamente materiale. L'LLM puo' proporla, ma l'accettazione richiede verifica contro la source e contro le altre proposition gia' estratte, soprattutto quando la tecnologia partecipa anche a produzione, servizio o transfer.
\end{ddtaprinciple}
%</PAGE-CONTENT:46>

\clearpage

%<PAGE-DESC:47>9 - realize: output BA atteso e discriminazioni differite
%<PAGE-CONTENT:47>
\PageHashFooter{\PageMDfive{47}}
\subsection{Output esatto atteso dall'estrazione assistita}
Il wrapper seguente e' un formato di authoring/review, non un nuovo costrutto BA. La parte di output BA usa soltanto la shape ammessa.

\begin{lstlisting}
REALIZE_ASSESSMENT
  decision = ACCEPT_REALIZE | NOT_A_REALIZE | INSUFFICIENT_EVIDENCE
  evidence = <source locator / concise supported wording>
  documentationGap = <abstract function missing | none>
  notSpecified = [<only unresolved material items>]

BA_OUTPUT
  BAReferent reuse/candidates:
    <canonical names only as needed>

  BAProposition <candidate-id>
    semanticOperatorKey = realize
    abstract    -> <BAReferent>
    realization -> <BAReferent>      # repeat for 1..*
    polarity = affirmative | negative
    scopedModifier = <condition|temporalScope>  # omit if absent
\end{lstlisting}

Se \texttt{decision != ACCEPT\_REALIZE}, \texttt{BA\_OUTPUT} non contiene una proposition \texttt{realize} accettabile. \texttt{documentationGap} appartiene al wrapper di review e non alla BA: serve soltanto a rendere visibile quando una tecnologia concreta e' documentata senza che sia governata la capability astratta che ci si aspetterebbe di poter sostituire o riesaminare.

\subsection{Discriminazioni con altri costrutti}
-
%</PAGE-CONTENT:47>

\clearpage

%<PAGE-DESC:48>9 - realize: rappresentazione Mermaid
%<PAGE-CONTENT:48>
\PageHashFooter{\PageMDfive{48}}
\subsection{Rappresentazione Mermaid di realize}

La projection usa una disposizione verticale simile alla relazione abstract/concrete: l'\texttt{abstract} e' sopra, la \texttt{realization} concreta sotto e una linea tratteggiata con triangolo vuoto punta verso l'abstract. La convenzione rende immediata la lettura ``il concreto realizza l'astratto'' senza trasformare la BA in un modello di classi.

\begin{tabularx}{\textwidth}{L{42mm}Y}
\toprule
\textbf{Parte di realize} & \textbf{Rappresentazione Mermaid} \\
\midrule
\texttt{abstract} & Nodo superiore, marcato graficamente come significato astratto. \\
\texttt{realization} & Nodo inferiore; ripetibile per \texttt{1..*}. \\
Relazione & Linea tratteggiata con triangolo vuoto rivolto verso l'abstract e label \texttt{realize}. \\
Proposition id / polarity & Metadati dell'intero pattern; non diventano nodi della gerarchia. \\
\bottomrule
\end{tabularx}

\begin{lstlisting}
classDiagram
direction BT

class ClinicalDescriptionGeneration {
  <<abstract>>
}

class Qwen2VL7BInstruct {
  <<realization>>
}

Qwen2VL7BInstruct ..|> ClinicalDescriptionGeneration : realize
\end{lstlisting}

\begin{ddtaanti}[title={Somiglianza grafica non equivale a subclassing}]
Il triangolo e la disposizione verticale sono una convenzione di projection per rendere visibile la relation abstract/concrete. Non autorizzano inferenze di ereditarieta', type hierarchy, interface conformance o object-oriented design.
\end{ddtaanti}
%</PAGE-CONTENT:48>

\clearpage

%<PAGE-DESC:49>9 - realize: esempio grafico Mermaid
%<PAGE-CONTENT:49>
\PageHashFooter{\PageMDfive{49}}
\subsection{Esempio grafico corrispondente}

Il diagramma applica la convenzione alla proposition DermaTriage \texttt{BAP-FR03-02-02}. \texttt{ClinicalDescriptionGeneration} resta il significato astratto; \texttt{Qwen2-VL-7B-Instruct} e' la realizzazione concreta della baseline. Il grafico non rappresenta input, output o transfer: questi appartengono ad altre proposition.

\vspace{8mm}
\begin{center}
\begin{tikzpicture}[
  abstractnode/.style={draw=DDTABlue,fill=DDTALightBlue,rounded corners=3pt,thick,align=center,minimum height=20mm,text width=56mm,inner sep=7pt,font=\sffamily\fontsize{8}{9.5}\selectfont},
  realizationnode/.style={draw=DDTABlue,fill=white,rounded corners=3pt,thick,align=center,minimum height=20mm,text width=56mm,inner sep=7pt,font=\sffamily\fontsize{8}{9.5}\selectfont},
  realizeflow/.style={draw=DDTABlue,very thick,dashed,-{Triangle[open,length=3.5mm,width=3.8mm]}}
]
\node[abstractnode] (a) {\textbf{ClinicalDescriptionGeneration}\\{\scriptsize\texttt{<<abstract>>}}};
\node[realizationnode,below=20mm of a] (r) {\textbf{Qwen2-VL-7B-Instruct}\\{\scriptsize\texttt{<<realization>>}}};

\draw[realizeflow] (r.north) -- node[right=3mm,font=\sffamily\scriptsize] {\texttt{realize}} (a.south);
\end{tikzpicture}
\end{center}

\vspace{5mm}
\begin{center}
\sffamily\small\textbf{BAP-FR03-02-02}\quad \texttt{realize}\quad \texttt{polarity = affirmative}
\end{center}
\vspace{4mm}

\begin{ddtacheck}[title={Lettura del grafico}]
La freccia si legge dal concreto verso l'astratto: \texttt{Qwen2-VL-7B-Instruct} realizza \texttt{ClinicalDescriptionGeneration}. Se una stessa proposition contenesse piu' realization governate, piu' rami tratteggiati convergerebbero sullo stesso abstract mantenendo un unico proposition id; il grafico non dichiarerebbe da solo se tali realization siano alternative o congiunte.
\end{ddtacheck}
%</PAGE-CONTENT:49>

\clearpage


%<PAGE-DESC:50>10 - constrain: definizione, smallest fact e criteri d'uso
%<PAGE-CONTENT:50>
\PageHashFooter{\PageMDfive{50}}
\section{Operatore \texttt{constrain}}
\CurrentTag

\subsection{Che cosa rappresenta}
\texttt{constrain} rappresenta una \textbf{restrizione governata, riusabile e interrogabile} applicata a un significato identificato. Il costrutto serve quando il progetto non si limita a nominare un dato, un comportamento o uno stato, ma ne governa anche un contratto, un dominio ammesso, una proprieta' richiesta o un'altra limitazione stabile che deve poter essere recuperata e analizzata separatamente.

\begin{ddtacore}[title={Smallest semantic fact}]
Esiste un target governato \texttt{T} e almeno una restriction \texttt{C} tale che la documentazione stabilisce che \texttt{T} sia soggetto a \texttt{C}. La restriction puo' essere un significato governato con identita' propria oppure una struttura locale controllata; in entrambi i casi il target resta esplicito.
\end{ddtacore}

\subsection{Quando usarlo}
Usare \texttt{constrain} quando la restriction deve restare distinguibile e interrogabile: per esempio un data/content contract, un allowed domain, una proprieta' obbligatoria, un requisito di confidenzialita' applicato a un delivery o un medium imposto a uno specifico comportamento di trasferimento.

\subsection{Quando non usarlo}
Non usare \texttt{constrain} per descrivere quale valore possiede una singola occorrenza runtime, quale ramo di una regola decisionale viene scelto, quale elemento viene selezionato, chi produce o trasferisce un dato, oppure quale meccanismo esegue materialmente il controllo. La restriction qualifica il target: non sostituisce il comportamento che produce, trasporta, seleziona o verifica.

\begin{ddtaopen}[title={Gate minimo}]
Prima di introdurre \texttt{constrain}, identificare il target esatto. Se non e' chiaro se la restriction appartenga al contenuto trasferito, al behavior di transfer, a un contratto riusabile o a un altro significato governato, non propagare il vincolo per convenienza: mantenere il gap esplicito e tornare alla documentazione quando necessario.
\end{ddtaopen}
%</PAGE-CONTENT:50>

\clearpage

%<PAGE-DESC:51>10 - constrain: firma canonica e forme di constraintValue
%<PAGE-CONTENT:51>
\PageHashFooter{\PageMDfive{51}}
\subsection{Firma canonica}
La firma top-level resta unica:

\begin{lstlisting}
constrain
  constraintTarget -> BAReferent [1]
  constraintValue  -> controlled value/structure [1..*]
  polarity         = affirmative | negative [1]
\end{lstlisting}

Gli scoped modifier \texttt{condition} e \texttt{temporalScope} compaiono soltanto quando la source limita realmente la restriction.

\begin{tabularx}{\textwidth}{L{34mm}L{20mm}Y}
\toprule
\textbf{Role} & \textbf{Card.} & \textbf{Significato} \\
\midrule
\texttt{constraintTarget} & 1 & BAReferent il cui significato e' ristretto dalla proposition. \\
\texttt{constraintValue} & 1..* & Restriction governata: puo' essere un BAReferent con identita' propria oppure una forma locale controllata. \\
\bottomrule
\end{tabularx}

\subsection{Due modi di esprimere la restriction}
\textbf{A. Constraint nominato / riusabile.} Quando il contratto o la restriction possiede identita' analitica indipendente:
\begin{lstlisting}
constraintValue -> <BAReferent>
\end{lstlisting}
Esempio: \texttt{ClinicalDescriptionContentContract}. Il nome non e' un'etichetta grafica: e' un BAReferent riusabile e change-addressable.

\textbf{B. Constraint locale.} Quando non serve una nuova identita' BA, la restriction resta inline. Puo' essere un dominio locale controllato, oppure una struttura di proprieta':
\begin{lstlisting}
constraintValue
  property   -> <controlled semantic key> [1]
  presence   -> required [0..1]
  vocabulary -> [<controlled typed local value | BAReferent> 1..*] [0..1]
\end{lstlisting}

Un allowed domain semplice puo' essere espresso direttamente come \texttt{constraintValue} locale, senza creare un BAReferent artificiale soltanto per ospitare l'elenco.
%</PAGE-CONTENT:51>

\clearpage

%<PAGE-DESC:52>10 - constrain: identity continuity, invarianti e forbidden inference
%<PAGE-CONTENT:52>
\PageHashFooter{\PageMDfive{52}}
\subsection{Continuita' di identita' e livello del constraint}
Il \texttt{constraintTarget} mantiene la propria identita' indipendentemente dalla restriction. Se \texttt{constraintValue} e' un BAReferent nominato, anche il constraint mantiene identita' separata e puo' a sua volta essere target di altre proposition.

\begin{lstlisting}
ClinicalDescription
  constrained by -> ClinicalDescriptionContentContract

ClinicalDescriptionContentContract
  constrained by -> structured local restrictions
\end{lstlisting}

Questa catena permette di distinguere \emph{a quale dato si applica il contratto} da \emph{come il contratto e' composto}. Una modifica della struttura interna del contratto non richiede di rinominare il dato o di duplicare tutte le proposition che vi fanno riferimento.

\subsection{Invarianti}
\begin{itemize}
  \item ogni proposition accettata possiede esattamente un \texttt{constraintTarget} e almeno un \texttt{constraintValue};
  \item un BAReferent usato come constraint nominato richiede normale giustificazione BA1 di identita' indipendente;
  \item \texttt{property} e i valori locali non diventano automaticamente BAReferent;
  \item \texttt{presence -> required} afferma la presenza della proprieta'; la sola \texttt{vocabulary} non la implica;
  \item piu' constraintValue restano nella stessa proposition soltanto se la stessa asserzione governata li applica allo stesso target e scope;
  \item una restriction sul content non si propaga automaticamente al behavior di transfer, e una restriction sul behavior non modifica automaticamente lo schema del content.
\end{itemize}

\subsection{Forbidden inference}
Dal solo \texttt{constrain} non inferire il valore runtime osservato, validazione riuscita, enforcement mechanism, algoritmo, ordine di valutazione, selection, mapping decisionale, produzione, transfer, ownership, responsibility, encryption effettivamente applicata o conformita' di una specifica occorrenza. Un allowed domain non dice quale valore verra' scelto; un data contract non dice da solo chi invia o riceve il dato.
%</PAGE-CONTENT:52>

\clearpage

%<PAGE-DESC:53>10 - constrain: domande di identificazione e procedura di estrazione
%<PAGE-CONTENT:53>
\PageHashFooter{\PageMDfive{53}}
\subsection{Domande da porre alla documentazione}
Le domande servono prima a localizzare la restriction e soltanto dopo a scegliere la sua forma.

\begin{enumerate}
  \item Qual e' il significato esatto che viene ristretto? Il target e' un dato/content, un behavior di transfer, uno stato, un contratto o altro BAReferent?
  \item La source governa realmente una restriction stabile e interrogabile oppure descrive soltanto un dettaglio locale senza utilita' analitica autonoma?
  \item La restriction possiede un nome/identita' riusabile nel progetto? Se si', esiste o serve un BAReferent dedicato?
  \item Se la restriction resta locale, e' un dominio diretto oppure una proprieta' strutturata?
  \item Per una proprieta' strutturata, qual e' il \texttt{property} key governato?
  \item La source richiede che la proprieta' sia presente? Usare \texttt{presence -> required} solo in questo caso.
  \item La source governa un insieme chiuso di valori ammessi? Inserirlo in \texttt{vocabulary} senza dedurre automaticamente la presenza.
  \item Il vincolo riguarda il dato stesso oppure il modo in cui quel dato viene trasferito? Non spostare il target tra content e transfer behavior.
  \item Piu' restriction appartengono davvero alla stessa asserzione e allo stesso scope oppure devono restare proposition separate?
  \item La restriction vale sempre oppure soltanto sotto condition/temporal scope governati?
  \item La source afferma o nega esplicitamente il constraint? Non dedurre polarity negativa dall'assenza di un valore o di una proprieta'.
  \item Se property, domain, scope o target non sono governati, lasciare \texttt{NOT SPECIFIED} / diagnosis invece di completarli per plausibilita'.
\end{enumerate}

\begin{ddtacheck}[title={Procedura minima}]
Identificare il target; determinare se la restriction richiede identita' propria oppure resta locale; compilare soltanto i dettagli effettivamente governati; verificare la separazione content/transfer behavior; infine riesaminare che la proposition non stia simulando una regola decisionale, una selection o un runtime fact.
\end{ddtacheck}
%</PAGE-CONTENT:53>

\clearpage

%<PAGE-DESC:54>10 - constrain: esempio generico e worked example DermaTriage corrente
%<PAGE-CONTENT:54>
\PageHashFooter{\PageMDfive{54}}
\subsection{Esempio generico minimo}
\begin{ddtaexample}[title={Documentazione inventata ma coerente}]
``Ogni \texttt{TelemetryRecord} trasferito tra i componenti MUST rispettare il \texttt{TelemetryRecordContract}. Il contratto richiede rappresentazione JSON e la presenza dei campi \texttt{eventId} e \texttt{timestamp}.''
\end{ddtaexample}

Assumendo identita' BA gia' giustificate:
\begin{lstlisting}
constrain
  constraintTarget -> TelemetryRecord
  constraintValue  -> TelemetryRecordContract
  polarity = affirmative

constrain
  constraintTarget -> TelemetryRecordContract
  constraintValue
    property   -> representation
    presence   -> required
    vocabulary -> [JSON]
  constraintValue
    property -> eventId
    presence -> required
  constraintValue
    property -> timestamp
    presence -> required
  polarity = affirmative
\end{lstlisting}

\subsection{Esempio corrente - DermaTriage Stage 2}
La documentazione corrente stabilisce che \texttt{ClinicalDescription} abbia un contratto di contenuto riusabile e che tale contratto sia testuale, composto da cinque parti richieste.

\begin{lstlisting}
BAP-FR03-02-03  constrain
  constraintTarget -> ClinicalDescription
  constraintValue  -> ClinicalDescriptionContentContract

BAP-FR03-02-04  constrain
  constraintTarget -> ClinicalDescriptionContentContract
  constraintValue -> representation required in [TEXTUAL]
  constraintValue -> partCount required in [5]
  constraintValue -> shape required
  constraintValue -> borders required
  constraintValue -> color required
  constraintValue -> texture required
  constraintValue -> suspiciousFeatures required
\end{lstlisting}

Il primo fatto identifica \emph{quale contratto governa il dato}; il secondo rende interrogabile \emph{come quel contratto e' composto}. I due livelli non vengono fusi.
%</PAGE-CONTENT:54>

\clearpage

%<PAGE-DESC:55>10 - constrain: prompt LLM per l'identificazione controllata
%<PAGE-CONTENT:55>
\PageHashFooter{\PageMDfive{55}}
\subsection{Prompt LLM di supporto all'authoring}
Il prompt deve preservare sia il target corretto sia il livello del constraint. In particolare non deve spostare una restriction dal dato al transfer behavior o viceversa soltanto per rendere il diagramma piu' comodo.

\begin{lstlisting}
TASK: identify governed BA constrain propositions.
Use ONLY the supplied documentation and existing BA registry.

A constrain proposition exists only if the source supports:
  constraintTarget [1], constraintValue [1..*].

Rules:
1. Identify the smallest governed reusable/queryable restriction fact.
2. Identify the exact target before extracting the restriction.
3. Reuse existing BAReferents when semantic identity is unchanged.
4. Use a named constraint BAReferent only when independent BA1 identity is
   justified; otherwise keep the restriction local.
5. For local structured constraints, use only governed property, presence and
   vocabulary information. Do not invent schema fields.
6. presence=required only when the source explicitly requires presence.
7. vocabulary does not imply presence.
8. Preserve the distinction between content constraints and transfer-behavior
   constraints. Do not move a restriction between them without source support.
9. Do not turn runtime values, validation results, decision mappings, selection,
   production or transfer semantics into constrain.
10. Keep multiple constraintValue entries together only when they share the same
    governed target and assertion scope.
11. Use polarity=negative only for an explicitly negated restriction assertion.
12. Add condition/temporalScope only when explicitly governed.
13. If target or every restriction value is missing, return INSUFFICIENT_EVIDENCE.
14. Do not invent operators, roles, states, BAReferents or metadata.

Return EXACTLY the authoring wrapper and BA output shape shown next.
\end{lstlisting}

\begin{ddtaprinciple}[title={Human review remains mandatory}]
L'LLM puo' riconoscere la forma del constraint, ma la review umana deve verificare soprattutto identity promotion, scope e placement: lo stesso vocabolo tecnico puo' qualificare il content, il behavior di transfer o un contratto nominato con effetti analitici diversi.
\end{ddtaprinciple}
%</PAGE-CONTENT:55>

\clearpage

%<PAGE-DESC:56>10 - constrain: output BA atteso e discriminazioni differite
%<PAGE-CONTENT:56>
\PageHashFooter{\PageMDfive{56}}
\subsection{Output esatto atteso dall'estrazione assistita}
Il wrapper distingue la decisione di authoring dalla BA. \texttt{constraintForm} serve alla review e non introduce un nuovo field BA.

\begin{lstlisting}
CONSTRAIN_ASSESSMENT
  decision = ACCEPT_CONSTRAIN | NOT_A_CONSTRAINT | INSUFFICIENT_EVIDENCE
  evidence = <source locator / concise supported wording>
  constraintForm = NAMED_REFERENT | LOCAL_DIRECT | LOCAL_STRUCTURED
  targetKind = CONTENT | TRANSFER_BEHAVIOR | OTHER
  documentationGap = <material unresolved meaning | none>
  notSpecified = [<only unresolved material items>]

BA_OUTPUT
  BAReferent reuse/candidates:
    <canonical names only as needed>

  BAProposition <candidate-id>
    semanticOperatorKey = constrain
    constraintTarget -> <BAReferent>
    constraintValue  -> <BAReferent | controlled local value/domain>
      # OR repeat local structured values as governed:
      property   -> <controlled semantic key>
      presence   -> required            # omit if absent
      vocabulary -> [<controlled values>] # omit if absent
    polarity = affirmative | negative
    scopedModifier = <condition|temporalScope>  # omit if absent
\end{lstlisting}

Se \texttt{decision != ACCEPT\_CONSTRAIN}, \texttt{BA\_OUTPUT} non contiene una proposition \texttt{constrain} accettabile. \texttt{constraintForm} e \texttt{targetKind} sono soltanto controlli di authoring per evitare che la scelta della rappresentazione alteri il significato governato.

\subsection{Discriminazioni con altri costrutti}
-
%</PAGE-CONTENT:56>

\clearpage

%<PAGE-DESC:57>10 - constrain: rappresentazione Mermaid
%<PAGE-CONTENT:57>
\PageHashFooter{\PageMDfive{57}}
\subsection{Rappresentazione Mermaid di constrain}

La projection mantiene visibile il \texttt{constraintTarget} e colloca la restriction \textbf{vicino al significato che qualifica}. Un constraint nominato usa un nodo dedicato; una restriction locale strutturata puo' essere resa come nota/card collegata al target. La linea tratteggiata etichettata \texttt{constrain} visualizza una BAProposition distinta e non modifica le frecce di data flow.

\begin{tabularx}{\textwidth}{L{42mm}Y}
\toprule
\textbf{Target del constraint} & \textbf{Placement nella view} \\
\midrule
Content / dato trasferito & Constraint accanto al nodo documento/content. Source e destination continuano a leggere lo stesso BAReferent informativo. \\
Transfer behavior & Constraint accanto al BAReferent del behavior/segmento di transfer, non al content. \\
Constraint nominato & Nodo riusabile collegato al target con linea tratteggiata \texttt{constrain}. \\
Constraint locale strutturato & Nota/card vicina al target, con le sole property/presence/vocabulary governate. \\
\bottomrule
\end{tabularx}

\begin{lstlisting}
flowchart LR
  S["ClinicalDescriptionGeneration<br/>source"]
  C@{ shape: doc, label: "ClinicalDescription<br/>content" }
  D["HistoricalCaseRetrieval<br/>destination"]
  K["ClinicalDescriptionContentContract<br/>named constraint"]
  F["TEXTUAL; 5 parts; required fields<br/>structured constraint"]

  S --> C --> D
  C -.->|constrain| K
  K -.->|constrain| F

  classDef endpoint fill:#EEF4F9,stroke:#275D86,color:#17324D
  classDef content fill:#FFFFFF,stroke:#275D86,color:#17324D
  classDef contract fill:#EAF5F3,stroke:#2A7F78,color:#17324D
  classDef detail fill:#F4F6F8,stroke:#5B6770,color:#17324D
  class S,D endpoint
  class C content
  class K contract
  class F detail
\end{lstlisting}

\begin{ddtaanti}[title={Placement non equivale a propagazione}]
Mostrare un data contract vicino a un transfer non significa che il transfer stesso sia il \texttt{constraintTarget}. Il target resta il content quando la restriction descrive forma o dominio del dato; diventa il transfer behavior soltanto quando la source governa una restriction sul delivery, per esempio medium o confidentiality del segmento.
\end{ddtaanti}
%</PAGE-CONTENT:57>

\clearpage

%<PAGE-DESC:58>10 - constrain: esempio grafico Mermaid composto con transfer
%<PAGE-CONTENT:58>
\PageHashFooter{\PageMDfive{58}}
\subsection{Esempio grafico corrispondente}

Il diagramma combina senza fonderle tre proposition: il \texttt{transfer} di \texttt{ClinicalDescription}, il constraint che associa il dato al \texttt{ClinicalDescriptionContentContract} e il constraint che descrive il contratto. Il data path resta \texttt{source -> content -> destination}; le linee tratteggiate appartengono a proposition \texttt{constrain} separate.

\vspace{6mm}
\begin{center}
\begin{tikzpicture}[
  endpoint/.style={draw=DDTABlue,fill=DDTALightBlue,rounded corners=3pt,thick,align=center,minimum height=15mm,text width=34mm,inner sep=4pt,font=\sffamily\fontsize{6.8}{7.8}\selectfont},
  contentdoc/.style={draw=DDTABlue,fill=white,thick,align=center,minimum height=22mm,text width=40mm,inner sep=7pt,font=\sffamily\fontsize{7}{8.2}\selectfont,
    path picture={\draw[DDTABlue,thick] ([xshift=-6mm]path picture bounding box.north east) -- ([xshift=-6mm,yshift=-6mm]path picture bounding box.north east) -- ([yshift=-6mm]path picture bounding box.north east); \draw[DDTABlue,thick] ([xshift=-6mm]path picture bounding box.north east) -- ([yshift=-6mm]path picture bounding box.north east);}},
  contract/.style={draw=DDTATeal,fill=DDTALightTeal,rounded corners=3pt,thick,align=center,minimum height=15mm,text width=42mm,inner sep=5pt,font=\sffamily\fontsize{6.8}{7.8}\selectfont},
  detail/.style={draw=DDTAGray,fill=DDTALightGray,rounded corners=2pt,thick,align=left,text width=54mm,inner sep=5pt,font=\sffamily\fontsize{6.2}{7.2}\selectfont},
  flow/.style={draw=DDTABlue,very thick,-{Stealth[length=3mm,width=2mm]}},
  constraint/.style={draw=DDTATeal,thick,dashed,-{Stealth[length=2.5mm,width=1.7mm]}}
]
\node[contentdoc] (c) at (0,0) {\textbf{ClinicalDescription}\\{\scriptsize\texttt{content}}};
\node[endpoint,left=18mm of c] (s) {\textbf{ClinicalDescriptionGeneration}\\{\scriptsize\texttt{source}}};
\node[endpoint,right=18mm of c] (d) {\textbf{HistoricalCaseRetrieval}\\{\scriptsize\texttt{destination}}};
\node[contract,below=16mm of c] (k) {\textbf{ClinicalDescriptionContentContract}\\{\scriptsize\texttt{named constraint}}};
\node[detail,below=11mm of k] (f) {\texttt{representation = TEXTUAL}\\\texttt{partCount = 5}\\\texttt{shape, borders, color, texture, suspiciousFeatures = required}};

\draw[flow] (s) -- (c);
\draw[flow] (c) -- (d);
\draw[constraint] (c.south) -- node[right=2mm,font=\sffamily\scriptsize]{\texttt{constrain}} (k.north);
\draw[constraint] (k.south) -- node[right=2mm,font=\sffamily\scriptsize]{\texttt{constrain}} (f.north);
\end{tikzpicture}
\end{center}

\vspace{3mm}
\begin{ddtacheck}[title={Lettura del grafico}]
Chi produce e chi riceve condividono lo stesso \texttt{ClinicalDescription}; il contratto rende visibile che cosa quel dato deve rispettare. Se una restriction riguardasse invece il delivery --- per esempio il medium o la confidentiality del trasferimento --- il nodo di constraint verrebbe collegato al BAReferent del transfer behavior e non a \texttt{ClinicalDescription}.
\end{ddtacheck}
%</PAGE-CONTENT:58>

\clearpage

\clearpage

%<PAGE-DESC:59>11 - condition: definizione, smallest fact e criteri d'uso
%<PAGE-CONTENT:59>
\PageHashFooter{\PageMDfive{59}}
\section{Scoped modifier \texttt{condition}}
\CurrentTag

\subsection{Che cosa rappresenta}
\texttt{condition} e' uno \textbf{scoped modifier BA2} che limita localmente l'applicabilita' di una BAProposition. Il fatto espresso dall'operatore e dai suoi participant non cambia identita': la condition stabilisce soltanto \emph{quando quella specifica asserzione si applica}. Non e' un nuovo top-level operator e non e' una decision rule abbreviata.

\begin{ddtacore}[title={Smallest semantic fact}]
Esiste una BAProposition governata \texttt{P} e una condizione locale \texttt{C} tale che \texttt{P} e' asserita soltanto quando \texttt{C} vale nello scope documentato. Eliminando \texttt{condition} si perderebbe il confine di applicabilita' della proposition, non il significato del suo operatore.
\end{ddtacore}

\subsection{Quando usarlo}
Usare \texttt{condition} quando la source governa un guard/applicability scope per una singola proposition e quel significato non richiede identita', provenance, participant o riuso autonomi. La condizione deve essere source-grounded; non puo' essere dedotta soltanto perche' un comportamento sembra plausibile in uno scenario.

\subsection{Quando non usarlo}
Non usare \texttt{condition} per codificare un mapping IF/THEN/ELSE che determina quale risultato viene prodotto, per una dependency/prerequisite, per un constraint riusabile o per nascondere un significato che deve diventare BAReferent/BAProposition autonomo.

\begin{ddtaopen}[title={Gate minimo}]
La domanda non e' ``c'e' un if nella frase?'', ma ``la documentazione limita quando questa stessa proposition vale?''. Se la condition seleziona risultati, introduce participant propri o deve essere riusata indipendentemente, fermarsi e riesaminare il costrutto.
\end{ddtaopen}
%</PAGE-CONTENT:59>

\clearpage

%<PAGE-DESC:60>11 - condition: forma canonica, scope locale e cardinalita
%<PAGE-CONTENT:60>
\PageHashFooter{\PageMDfive{60}}
\subsection{Forma canonica nella BAProposition}
La grammatica BA2 preservata ammette \texttt{scopedModifier [0..*]} e, nel lower bound corrente, soltanto i kind \texttt{condition} e \texttt{temporalScope}. R7 rende esplicita \texttt{condition} con la seguente forma di authoring:

\begin{lstlisting}
BAProposition
  semanticOperatorKey = <operator>
  participation ...
  polarity = affirmative | negative

  scopedModifier
    condition -> <source-grounded local applicability statement>
\end{lstlisting}

La parte destra di \texttt{condition} e' contenuto locale controllato della proposition, non un BAReferent implicito. La serializzazione deve preservare la formulazione canonica scelta durante la review senza inventare property, boolean o stati non governati.

\begin{tabularx}{\textwidth}{L{38mm}L{22mm}Y}
\toprule
\textbf{Elemento} & \textbf{Card.} & \textbf{Significato} \\
\midrule
\texttt{scopedModifier} & 0..* & Modifier locali ammessi dalla BAProposition; il lower bound corrente contiene \texttt{condition} e \texttt{temporalScope}. \\
\texttt{condition} & per instance 1 & Guard/applicability statement governato per la proposition che lo contiene. \\
\bottomrule
\end{tabularx}

\begin{ddtaprinciple}[title={Nessun DSL logico generale}]
La forma locale non introduce automaticamente \texttt{AND}, \texttt{OR}, truth table, property model o evaluation order. Se la source governa una costruzione decisionale che seleziona un risultato, il candidato da riesaminare e' \texttt{decisionRule}, non una condition sempre piu' complessa.
\end{ddtaprinciple}
%</PAGE-CONTENT:60>

\clearpage

%<PAGE-DESC:61>11 - condition: identity boundary, invarianti e forbidden inference
%<PAGE-CONTENT:61>
\PageHashFooter{\PageMDfive{61}}
\subsection{Locality / identity boundary}
Una condition resta embedded soltanto finche' modifica localmente una proposition e non richiede identita' analitica indipendente. BA2 R2/R3 preserva il promotion test: se il significato viene riusato, separatamente revisionato, diventa participant-bearing o deve essere target di altre asserzioni, si applicano i normali criteri BA1 invece di nasconderlo nel modifier.

\begin{lstlisting}
LOCAL
  produce(...)
  condition -> LesionImage unavailable

PROMOTION PRESSURE
  NoImageCondition must be reused / constrained / traced independently
  -> evaluate BAReferent / BAProposition identity under BA1
\end{lstlisting}

\subsection{Invarianti}
\begin{itemize}
  \item la condition qualifica l'intera BAProposition, non un singolo edge scelto dal renderer;
  \item l'operatore, i roleKey e i BAReferent della proposition mantengono la stessa identita';
  \item una condition locale non crea automaticamente un nuovo BAReferent;
  \item l'assenza della condition opposta non autorizza a inventare un ramo \texttt{ELSE};
  \item l'ordine testuale di piu' facts o view non implica evaluation order;
  \item la condition deve essere governata dalla source allo stesso livello documentale che giustifica la proposition o da evidence tracciabile compatibile con tale scope.
\end{itemize}

\subsection{Forbidden inference}
Da \texttt{condition} non inferire detection logic, boolean property, stato runtime persistente, branch complementare, fallback, precedence, sequencing, retry/failure semantics, dependency, constraint riusabile o mapping di risultato. ``Image unavailable'' non autorizza da solo a inventare \texttt{LesionImage.available = FALSE}.
%</PAGE-CONTENT:61>

\clearpage

%<PAGE-DESC:62>11 - condition: domande di identificazione e procedura di estrazione
%<PAGE-CONTENT:62>
\PageHashFooter{\PageMDfive{62}}
\subsection{Domande da porre alla documentazione}
\begin{enumerate}
  \item Quale BAProposition e' limitata dalla condizione? Il modifier deve avere un owner propositionale preciso.
  \item La source sta dicendo \emph{quando la proposition vale} oppure \emph{come una condizione determina quale risultato viene scelto}?
  \item La condition e' esplicita/governata oppure la stiamo deducendo da un percorso ritenuto probabile?
  \item Il significato puo' restare locale senza BAReferent autonomo?
  \item E' riusato, separatamente revisionato, participant-bearing o target di altre proposition? Se si', applicare l'identity test BA1.
  \item La formulazione richiederebbe di inventare property, boolean, enum o stato? Se si', non completarli per plausibilita'.
  \item Esiste davvero un ramo opposto governato? Una condition non implica automaticamente \texttt{ELSE} o il complemento logico.
  \item La source governa una dependency/prerequisite? In tal caso verificare \texttt{dependOn}, non \texttt{condition}.
  \item La source governa una restriction riusabile? In tal caso verificare \texttt{constrain}.
  \item La condition e' temporale? Se descrive before/after/during/until, verificare \texttt{temporalScope} invece di condition.
\end{enumerate}

\begin{ddtacheck}[title={Procedura minima}]
Estrarre prima la BAProposition senza modifier; verificare che il fatto sia semanticamente completo; isolare poi il guard governato; applicare locality/identity test; aggiungere \texttt{condition} soltanto se cambia l'applicabilita' della stessa proposition. Infine controllare che non stia simulando decisionRule, dependOn o constrain.
\end{ddtacheck}
%</PAGE-CONTENT:62>

\clearpage

%<PAGE-DESC:63>11 - condition: esempio generico e worked example DermaTriage corrente
%<PAGE-CONTENT:63>
\PageHashFooter{\PageMDfive{63}}
\subsection{Esempio generico minimo}
\begin{ddtaexample}[title={Documentazione inventata ma coerente}]
``Quando la modalita' di manutenzione e' attiva, \texttt{SystemMonitor} produce un \texttt{MaintenanceAlert} a partire da \texttt{SensorStatus}.''
\end{ddtaexample}

Assumendo i BAReferent gia' giustificati:
\begin{lstlisting}
produce
  actor  -> SystemMonitor
  input  -> SensorStatus
  result -> MaintenanceAlert
  polarity = affirmative
  scopedModifier
    condition -> MaintenanceMode active
\end{lstlisting}

La condition dice quando questa produzione si applica. Non afferma che \texttt{MaintenanceMode} sia un boolean, non crea un ramo per la modalita' inattiva e non descrive come viene rilevato lo stato.

\subsection{Esempio corrente - DermaTriage symptom-only}
La documentazione R25 stabilisce: ``Quando l'immagine della lesione non e' disponibile, DermaTriage determina l'urgenza del caso utilizzando le informazioni sintomatologiche disponibili associate allo stesso caso.''

\begin{lstlisting}
produce
  actor  -> DermaTriage
  input  -> SymptomInformation
  result -> CaseUrgency
  polarity = affirmative
  scopedModifier
    condition -> LesionImage unavailable
\end{lstlisting}

Non serve materializzare \texttt{NoImageAvailableCondition} come BAReferent finche' il significato resta locale. Il percorso image-based con immagine disponibile deve essere supportato da proposition/documentazione proprie; non viene inventato come \texttt{ELSE} di questa proposition.
%</PAGE-CONTENT:63>

\clearpage

%<PAGE-DESC:64>11 - condition: prompt LLM per l'identificazione controllata
%<PAGE-CONTENT:64>
\PageHashFooter{\PageMDfive{64}}
\subsection{Prompt LLM di supporto all'authoring}
\begin{lstlisting}
TASK: identify governed proposition-local condition modifiers.
Use ONLY the supplied documentation, accepted BA and current BA authority.

A condition modifier exists only when the source limits when an otherwise
complete BAProposition applies.

Rules:
1. Identify the complete BAProposition first: operator, roles and referents.
2. Extract condition only if the source explicitly governs the applicability guard.
3. The condition qualifies the whole proposition, not an arbitrary visual edge.
4. Keep the condition local when it has no independent identity/reuse.
5. Do not create a BAReferent only to name a one-use guard.
6. If the condition is reused, separately reviewed, participant-bearing or targeted
   by other propositions, stop and request BA1 identity review.
7. Do not invent boolean properties, enums, states or opposite conditions.
8. Do not infer ELSE, fallback or complementary behavior from one condition.
9. Do not use condition for IF/THEN/ELSE result construction; review decisionRule.
10. Do not use condition for prerequisite/dependency; review dependOn.
11. Do not use condition for reusable restrictions; review constrain.
12. Do not infer temporal order; temporal applicability belongs to temporalScope.
13. If the guard is not source-grounded, return INSUFFICIENT_EVIDENCE.
14. Do not invent operators, roles, BAReferents or metadata.

Return EXACTLY the authoring wrapper and BA output shape shown next.
\end{lstlisting}

\begin{ddtaprinciple}[title={Human review remains mandatory}]
La review umana deve verificare soprattutto il confine fra applicabilita' locale e semantica autonoma. La forma linguistica ``quando/se'' e' soltanto un indizio: non decide automaticamente condition, decisionRule o dependOn.
\end{ddtaprinciple}
%</PAGE-CONTENT:64>

\clearpage

%<PAGE-DESC:65>11 - condition: output BA atteso e discriminazioni
%<PAGE-CONTENT:65>
\PageHashFooter{\PageMDfive{65}}
\subsection{Output esatto atteso dall'estrazione assistita}
\begin{lstlisting}
CONDITION_ASSESSMENT
  decision = ACCEPT_CONDITION | NOT_A_CONDITION | INSUFFICIENT_EVIDENCE
  evidence = <source locator / concise supported wording>
  ownerProposition = <candidate-id / identified governed fact>
  identityPressure = NONE | REVIEW_BA1
  notSpecified = [<only unresolved material items>]

BA_OUTPUT
  BAReferent reuse/candidates:
    <canonical names only as needed>

  BAProposition <candidate-id>
    semanticOperatorKey = <existing operator>
    <operator-scoped roles and terms>
    polarity = affirmative | negative
    scopedModifier
      condition -> <source-grounded local applicability statement>
\end{lstlisting}

\texttt{ownerProposition} e \texttt{identityPressure} appartengono al wrapper di review, non alla BA.

\subsection{Discriminazioni con i costrutti vicini}
\begin{tabularx}{\textwidth}{L{32mm}Y}
\toprule
\textbf{Confine} & \textbf{Discriminatore} \\
\midrule
condition / decisionRule & condition cambia quando la proposition si applica; decisionRule governa come condizioni/input determinano un risultato. \\
condition / constrain & condition e' locale all'asserzione; constrain e' una restriction riusabile/interrogabile sul proprio target. \\
condition / dependOn & condition e' guard di applicabilita'; dependOn afferma prerequisite/dependency e supporta impact propagation. \\
condition / BAReferent & il modifier resta locale; se il significato richiede identita'/riuso autonomi, applicare BA1. \\
condition / temporalScope & condition governa un guard non temporale; before/after/during/until appartengono al temporal scope. \\
\bottomrule
\end{tabularx}
%</PAGE-CONTENT:65>

\clearpage

%<PAGE-DESC:66>11 - condition: rappresentazione Mermaid e semantic determinism
%<PAGE-CONTENT:66>
\PageHashFooter{\PageMDfive{66}}
\subsection{Rappresentazione Mermaid di condition}
La \texttt{condition} qualifica l'intera BAProposition. La projection canonica la rende come un \textbf{tag di condition} colorato, posto sopra il ramo governato. Il tag non e' un BAReferent, non e' un input dell'operatore e non e' un nodo decisionale: visualizza soltanto lo \texttt{scopedModifier.condition} della proposition owner o del blocco di vista selezionato per quel ramo.

\begin{tabularx}{\textwidth}{L{44mm}Y}
\toprule
\textbf{Semantica BA} & \textbf{Regola di projection} \\
\midrule
\texttt{condition} modifier & Un tag ellittico colorato con label canonica \texttt{<condition>}. \\
Applicabilita' & Il tag viene posto in testa al ramo governato; la sua posizione rende visibile l'applicabilita' locale senza introdurre dataflow o nuovi role BA. \\
Owner block & Nelle viste composte il ramo puo' essere rappresentato con \emph{un solo blocco owner} per lato, quando l'obiettivo e' mostrare la separazione di condition e non l'espansione completa di tutte le proposition sottostanti. \\
Due condition distinte & Producono due tag distinti; nessun ramo viene inventato come complemento dell'altro. \\
Decision rule & Il rombo decisionale resta riservato alla projection di \texttt{decisionRule}; \texttt{condition} non usa il rombo e non introduce automaticamente una forma \texttt{IF/ELSE}. \\
\bottomrule
\end{tabularx}

\begin{lstlisting}
flowchart LR
  C1([Condition A])
  B1["Owner block A"]

  C2([Condition B])
  B2["Owner block B"]

  C1 --> B1
  C2 --> B2
\end{lstlisting}

\begin{ddtacore}[title={Semantic determinism della projection}]
Con la stessa BA accettata e lo stesso projection profile devono restare invariati: numero dei tag di condition, shape ellittica, label canonica della condition e ramo owner. La posizione geometrica dei nodi puo' variare: in questa fase non richiediamo layout determinism. Se la vista usa un solo blocco owner per ramo, anche quella scelta deve essere esplicita e coerente con lo scope della view.
\end{ddtacore}
%</PAGE-CONTENT:66>

\clearpage

%<PAGE-DESC:67>11 - condition: esempio grafico Mermaid composto con produce
%<PAGE-CONTENT:67>
\PageHashFooter{\PageMDfive{67}}
\subsection{Esempio grafico corrispondente}
Il diagramma mostra sullo stesso foglio due rami con \texttt{condition} distinte. I tag colorati \texttt{LesionImage available} e \texttt{LesionImage unavailable} stanno sopra i rispettivi rami. Per mantenere la vista leggibile, ogni lato della condition e' rappresentato con \textbf{un solo blocco owner} invece di espandere tutto cio' che sarebbe ulteriormente derivabile dalla documentazione.

\vspace{5mm}
\begin{center}
\begin{tikzpicture}[
  owner/.style={draw=DDTABlue,fill=white,rounded corners=3pt,thick,align=center,minimum height=18mm,text width=38mm,inner sep=6pt,font=\sffamily\fontsize{7}{8.2}\selectfont},
  actor/.style={draw=DDTABlue,fill=DDTALightBlue,rounded corners=3pt,thick,align=center,minimum height=16mm,text width=34mm,inner sep=5pt,font=\sffamily\fontsize{7}{8.2}\selectfont},
  conditiontag/.style={ellipse,draw=DDTAAmber,fill=DDTALightAmber,thick,align=center,minimum width=34mm,minimum height=11mm,inner sep=3pt,font=\sffamily\fontsize{7}{8}\selectfont},
  conn/.style={draw=DDTABlue,thick,-{Stealth[length=2.7mm,width=1.8mm]}}
]
\node[actor] (eval) at (0,4.6) {\textbf{Dermatological}\\\textbf{TriageEvaluation}};

\node[conditiontag] (c1) at (-4.4,3.0) {\texttt{LesionImage available}};
\node[conditiontag] (c2) at (4.4,3.0) {\texttt{LesionImage unavailable}};

\draw[conn] (eval.south west) .. controls (-1.9,3.95) and (-3.2,3.45) .. (c1.north);
\draw[conn] (eval.south east) .. controls (1.9,3.95) and (3.2,3.45) .. (c2.north);

\node[owner] (b1) at (-4.4,0.8) {\textbf{Image-based}\\\textbf{triage path}\\{\scriptsize branch owner block}};
\node[owner] (b2) at (4.4,0.8) {\textbf{Symptom-based}\\\textbf{urgency path}\\{\scriptsize branch owner block}};

\draw[conn] (c1.south) -- (b1.north);
\draw[conn] (c2.south) -- (b2.north);
\end{tikzpicture}
\end{center}

\vspace{2mm}
\begin{center}
\sffamily\small
\texttt{condition = LesionImage available}\qquad\texttt{condition = LesionImage unavailable}
\end{center}

\begin{ddtacheck}[title={Lettura del grafico}]
Il tag colorato indica soltanto la condition del ramo. Non e' un nodo decisionale e non implica automaticamente un ramo complementare. In questa vista ciascun lato e' compresso in un solo blocco owner: la projection evidenzia il fatto che due condition diverse governano due rami diversi, ma non espande automaticamente l'intera struttura interna dei rami.
\end{ddtacheck}
%</PAGE-CONTENT:67>

\clearpage
%<PAGE-DESC:68>12 - decisionRule: definizione, smallest fact e criteri d'uso
%<PAGE-CONTENT:68>
\PageHashFooter{\PageMDfive{68}}
\section{Operator structure \texttt{decisionRule}}
\CandidateTag

\subsection{Che cosa rappresenta}
\texttt{decisionRule} e' qui una \textbf{candidate operator structure} locale a una BAProposition \texttt{produce}. Rappresenta il caso in cui la source non governa soltanto che un actor produce un result da uno o piu' input, ma governa anche \emph{come condizioni o valori degli input selezionano il valore del result}. La candidate non rinomina il costrutto: riesamina soltanto la collocazione dello storico \texttt{decisionRule}. Actor, input e result restano quelli dell'owner \texttt{produce}; la struttura decisionale viene annidata.

\begin{ddtacore}[title={Smallest semantic fact}]
Esiste una proposition \texttt{produce} gia' completa e una regola governata \texttt{R} che determina quale valore del suo result viene selezionato a partire dagli input/condizioni documentati. Eliminando \texttt{decisionRule} si conserva il fatto che il result viene prodotto, ma si perde il mapping operativo governato che seleziona il valore del result.
\end{ddtacore}

\subsection{Quando usarla}
Usare la candidate quando la documentazione esplicita un mapping del tipo ``in queste condizioni, il result assume questo valore'' e actor/input/result coincidono semanticamente con quelli di una production gia' identificata. La regola deve essere governata: non basta che un mapping sia plausibile o implementato nel codice.

\subsection{Quando non usarla}
Non usarla per dire soltanto quando la proposition si applica: quello resta \texttt{scopedModifier.condition}. Non usarla per dichiarare soltanto il dominio ammesso del result: quello resta \texttt{constrain}. Non usarla per dependency, sequencing o fallback non governati.

\begin{ddtaopen}[title={Delta rispetto all'authority corrente}]
R3 mantiene \texttt{decisionRule} come top-level operator con propri role \texttt{actor/input/result}. R7 non cambia tale authority: testa se lo stesso \texttt{decisionRule} possa essere rilocato come \texttt{operatorStructure} dell'owner \texttt{produce}, evitando la duplicazione di actor/input/result. La candidate richiede regression e promotion esplicita.
\end{ddtaopen}
%</PAGE-CONTENT:68>

\clearpage

%<PAGE-DESC:69>12 - decisionRule: forma canonica candidate e cardinalita
%<PAGE-CONTENT:69>
\PageHashFooter{\PageMDfive{69}}
\subsection{Forma canonica candidate dentro \texttt{produce}}
La forma di rebuild mantiene la proposition \texttt{produce} come owner del fatto produttivo e aggiunge al massimo una struttura locale di selezione del risultato:

\begin{lstlisting}
BAProposition
  semanticOperatorKey = produce
  participation
    actor  -> <BAReferent> [1]
    input  -> <BAReferent> [0..*]
    result -> <BAReferent> [1..*]
  polarity = affirmative | negative
  scopedModifier [0..*]
    condition -> <local applicability guard>
  operatorStructure [0..1]
    decisionRule
      branch [1..*]
        when -> <source-grounded selection condition> [1]
        resultAssignment [1..*]
          target -> <result BAReferent>
          value  -> <controlled local value | BAReferent>
\end{lstlisting}

\begin{tabularx}{\textwidth}{L{39mm}L{22mm}Y}
\toprule
\textbf{Elemento} & \textbf{Card.} & \textbf{Significato} \\
\midrule
\texttt{operatorStructure} & 0..1 & Struttura locale dell'owner proposition; non crea da sola una seconda BAProposition. \\
\texttt{decisionRule} & per owner 1 & Mapping governato che seleziona il valore del result. \\
\texttt{branch} & 1..* & Branch esplicitamente governati dalla source. \\
\texttt{when} & per branch 1 & Condizione di selezione source-grounded; puo' preservare una residual condition esplicita senza inventare un ELSE globale. \\
\texttt{resultAssignment} & 1..* & Assegnazione governata a un result gia' partecipante all'owner \texttt{produce}. \\
\bottomrule
\end{tabularx}

\begin{ddtaopen}[title={Vocabolario delle comparison ancora sotto pressure}]
L'authority corrente di \texttt{decisionRule.comparison} ammette \texttt{equals|notEquals}. DermaTriage governa esplicitamente \texttt{confidence > 0.85} e quindi mette sotto pressure l'aggiunta candidate di \texttt{greaterThan}. Non generalizzare automaticamente a \texttt{>=}, \texttt{<}, range o aritmetica senza nuovi casi governati.
\end{ddtaopen}
%</PAGE-CONTENT:69>

\clearpage

%<PAGE-DESC:70>12 - decisionRule: identity boundary, invarianti e forbidden inference
%<PAGE-CONTENT:70>
\PageHashFooter{\PageMDfive{70}}
\subsection{Locality / identity boundary}
Nella candidate, il \texttt{decisionRule} annidato condivide l'identita' della BAProposition \texttt{produce} che qualifica. Non viene creato un secondo fact con gli stessi actor/input/result soltanto per ospitare il mapping. Questo e' precisamente il delta da verificare contro l'attuale \texttt{decisionRule} top-level dell'authority.

\begin{lstlisting}
OWNER FACT
  produce(actor, input(s), result)

LOCAL STRUCTURE
  operatorStructure.decisionRule
    -> how the governed result value is selected
\end{lstlisting}

Se la regola deve avere provenance, review/change lifecycle, riuso o targetability realmente indipendenti dall'owner \texttt{produce}, la rilocazione locale non va forzata: il caso riapre l'identity/BA2 design e puo' falsificare questa candidate.

\subsection{Invarianti}
\begin{itemize}
  \item actor, input e result sono quelli della proposition \texttt{produce}; la struttura non li duplica;
  \item \texttt{condition} e \texttt{decisionRule} possono coesistere: la prima governa quando la production si applica, la seconda quale valore del result viene selezionato una volta applicabile;
  \item ogni branch deve essere source-grounded; un branch mancante non viene inventato come \texttt{ELSE};
  \item una residual condition come ``gli altri casi HIGH'' resta limitata allo scope governato e non diventa un default globale;
  \item l'ordine dei branch non implica precedence a meno che la source non governi un ordine necessario;
  \item il mapping non modifica il dominio ammesso del result se tale dominio e' governato separatamente da \texttt{constrain}.
\end{itemize}

\subsection{Forbidden inference}
Non inferire completezza del mapping, default, precedence, short-circuit, gestione degli input mancanti, coercion dei tipi, range non documentati, error/failure handling o nuovi comparison key. Non trasformare valori locali come \texttt{P1} in BAReferent se non superano l'identity test BA1.
%</PAGE-CONTENT:70>

\clearpage

%<PAGE-DESC:71>12 - decisionRule: domande di identificazione e procedura di estrazione
%<PAGE-CONTENT:71>
\PageHashFooter{\PageMDfive{71}}
\subsection{Domande da porre alla documentazione}
\begin{enumerate}
  \item Esiste gia' un fatto \texttt{produce} completo con actor, input e result governati?
  \item La source sta dicendo \emph{quando la production vale} o \emph{quale valore del result viene selezionato}?
  \item Le condizioni del mapping sono esplicite e governate, oppure servirebbe inventare property o soglie?
  \item I branch producono valori dello stesso result della proposition owner?
  \item Il dominio dei valori e' governato separatamente e richiede anche un \texttt{constrain}?
  \item Esistono branch residuali espliciti (``altri casi'', ``otherwise'')? Qual e' il loro scope esatto?
  \item I branch sono semanticamente disgiunti oppure servirebbe precedence/evaluation order non documentato?
  \item Le comparison richieste appartengono al vocabolario corrente? Se no, registrare pressure invece di inventare operatori di confronto.
  \item La regola deve essere riusata, revisionata o tracciata indipendentemente dall'owner \texttt{produce}? Se si', fermare la rilocazione locale e riaprire identity/BA2 review.
  \item La source governa davvero un mapping di risultato oppure soltanto una dependency, restriction o classification?
\end{enumerate}

\begin{ddtacheck}[title={Procedura minima}]
Estrarre prima \texttt{produce} senza rule; stabilizzare actor/input/result; separare eventuale \texttt{condition} di applicabilita'; isolare poi soltanto i branch che selezionano il result; preservare i valori e le condizioni governate; verificare branch residuali e comparison vocabulary; infine confrontare la collocazione candidate annidata con l'attuale \texttt{decisionRule} top-level per assicurare che nessuna identita' o provenance necessaria venga persa.
\end{ddtacheck}
%</PAGE-CONTENT:71>

\clearpage

%<PAGE-DESC:72>12 - decisionRule: esempio generico e pressure case DermaTriage
%<PAGE-CONTENT:72>
\PageHashFooter{\PageMDfive{72}}
\subsection{Esempio generico minimo}
\begin{ddtaexample}[title={Documentazione inventata ma coerente}]
``\texttt{QueueClassifier} produce \texttt{QueueClass} da \texttt{RequestClassification}. Se \texttt{classKind} e' \texttt{EXPRESS}, il risultato e' \texttt{FAST}; se \texttt{classKind} e' \texttt{STANDARD}, il risultato e' \texttt{NORMAL}.''
\end{ddtaexample}

\begin{lstlisting}
produce
  actor  -> QueueClassifier
  input  -> RequestClassification
  result -> QueueClass
  polarity = affirmative
  operatorStructure
    decisionRule
      branch
        when -> RequestClassification.classKind equals EXPRESS
        resultAssignment -> QueueClass = FAST
      branch
        when -> RequestClassification.classKind equals STANDARD
        resultAssignment -> QueueClass = NORMAL
\end{lstlisting}

L'assenza di altri branch non autorizza a inventare un default.

\subsection{Pressure case corrente - DermaTriage}
\texttt{FR-MR01-04-01} governa quattro selezioni: \texttt{HIGH + confidence > 0.85 -> P1}; gli altri casi \texttt{HIGH -> P2}; \texttt{MEDIUM -> P3}; \texttt{LOW -> P4}. La seguente e' una \emph{candidate authoring view}, non BA gia' accettata:

\begin{lstlisting}
produce
  actor  -> OperationalPriorityDerivation   [candidate]
  input  -> CaseUrgency
  input  -> TriageConfidence
  result -> OperationalTriagePriority
  scopedModifier
    condition -> LesionImage available
  operatorStructure
    decisionRule
      HIGH + confidence > 0.85 -> P1
      HIGH, other applicable cases -> P2
      MEDIUM -> P3
      LOW -> P4
\end{lstlisting}

Il caso forza due review separate: \texttt{greaterThan} nel comparison vocabulary e una forma multi-branch che preservi ``altri casi HIGH'' senza trasformarla in ELSE globale.
%</PAGE-CONTENT:72>

\clearpage

%<PAGE-DESC:73>12 - decisionRule: prompt LLM per l'identificazione controllata
%<PAGE-CONTENT:73>
\PageHashFooter{\PageMDfive{73}}
\subsection{Prompt LLM di supporto all'authoring}
\begin{lstlisting}
TASK: identify governed result-selection logic for an existing produce proposition.
Use ONLY supplied documentation, accepted BA and current BA authority.
This R7 nesting of decisionRule is CANDIDATE and does not replace current top-level decisionRule authority.

Rules:
1. Identify the complete produce proposition first: actor, input(s), result(s), polarity.
2. Separate proposition applicability from result selection.
3. Put applicability in scopedModifier.condition, never in operatorStructure.decisionRule.
4. Identify nested decisionRule only if the source governs how conditions/input
   values select a value of an existing produce result.
5. Reuse the produce actor/input/result. Do not duplicate them as a second fact.
6. Extract only source-governed branches and assignments.
7. Do not invent ELSE/default/completeness/precedence/evaluation order.
8. Preserve explicit residual wording such as "other HIGH cases" in its exact scope.
9. Use only admitted comparison forms; if the source requires a new comparison
   (for example greaterThan), return VOCABULARY_PRESSURE rather than inventing it.
10. Do not turn allowed-domain semantics into result selection; review constrain.
11. If the rule needs independent identity/reuse/provenance/change lifecycle,
    return IDENTITY_PRESSURE and stop local nesting.
12. Do not invent BAReferents, properties, typed values or metadata.

Return EXACTLY the review wrapper and candidate BA shape shown next.
\end{lstlisting}

\begin{ddtaprinciple}[title={Human review remains mandatory}]
La review umana deve verificare due confini distinti: \texttt{condition} contro result selection e local structure contro independent proposition identity. La presenza linguistica di ``if'' non decide nessuno dei due confini.
\end{ddtaprinciple}
%</PAGE-CONTENT:73>

\clearpage

%<PAGE-DESC:74>12 - decisionRule: output candidate e discriminazioni
%<PAGE-CONTENT:74>
\PageHashFooter{\PageMDfive{74}}
\subsection{Output esatto atteso dall'estrazione assistita}
\begin{lstlisting}
RESULT_SELECTION_ASSESSMENT
  decision = ACCEPT_CANDIDATE | NOT_RESULT_SELECTION | INSUFFICIENT_EVIDENCE
  evidence = <source locator / concise supported wording>
  ownerProduce = <candidate-id / identified governed fact>
  authorityStatus = CANDIDATE_RELOCATION_OF_DECISION_RULE
  comparisonPressure = NONE | VOCABULARY_PRESSURE
  identityPressure = NONE | REVIEW_BA1_BA2
  notSpecified = [<only unresolved material items>]

BA_OUTPUT_CANDIDATE
  BAProposition <owner-id>
    semanticOperatorKey = produce
    <produce roles and terms>
    polarity = affirmative | negative
    scopedModifier [only if independently source-grounded]
      condition -> <applicability guard>
    operatorStructure
      decisionRule
        <source-grounded branches and result assignments>
\end{lstlisting}

I campi di assessment appartengono al wrapper di review, non alla BA.

\subsection{Discriminazioni con i costrutti vicini}
\begin{tabularx}{\textwidth}{L{37mm}Y}
\toprule
\textbf{Confine} & \textbf{Discriminatore} \\
\midrule
decisionRule / condition & la rule seleziona il valore del result; condition cambia quando l'intera proposition si applica. \\
decisionRule / constrain & la rule governa la selezione runtime/operativa; constrain governa un dominio o restriction senza dire come un valore viene scelto. \\
decisionRule / produce & produce resta il fatto owner; la rule aggiunge dettaglio di selezione senza creare un secondo actor/input/result. \\
nested decisionRule / top-level decisionRule & e' lo stesso costrutto canonico: l'authority corrente lo rappresenta come operator top-level; la candidate R7 ne testa soltanto la rilocazione come \texttt{operatorStructure} dell'owner \texttt{produce}. \\
decisionRule / dependOn & la rule seleziona un risultato; dependOn dichiara prerequisite/dependency e non simula branch IF/THEN. \\
\bottomrule
\end{tabularx}
%</PAGE-CONTENT:74>

\clearpage

%<PAGE-DESC:75>12 - decisionRule: rappresentazione Mermaid e visibility profile
%<PAGE-CONTENT:75>
\PageHashFooter{\PageMDfive{75}}
\subsection{Rappresentazione Mermaid della operator structure}
Quando una view sceglie di renderizzare \texttt{decisionRule}, la struttura appare come \textbf{annotazione grigia a bordo tratteggiato} associata alla proposition \texttt{produce}. Non e' un BAReferent, non e' un secondo actor e non crea un nuovo dataflow edge.

\begin{tabularx}{\textwidth}{L{42mm}Y}
\toprule
\textbf{Semantica BA} & \textbf{Regola di projection candidate} \\
\midrule
\texttt{produce} & Mantiene il normale pattern blu \texttt{input(s) -> actor -> result(s)}. \\
\texttt{condition} & Se inclusa dalla view, resta il tag ambra sopra il ramo/proposition owner. \\
\texttt{decisionRule} & Box grigio chiaro, bordo tratteggiato, collegamento di qualificazione senza arrowhead di dataflow. \\
Visibility & Il projection profile puo' omettere completamente il box quando la decision logic non interessa alla view. L'omissione non crea uno stato BA ``collapsed''. \\
Dettaglio & Se renderizzata, la view puo' mostrare i branch governati necessari al proprio scope; non puo' cambiarne significato o inventarne altri. \\
\bottomrule
\end{tabularx}

\begin{lstlisting}
flowchart LR
  I1["Input A"] --> A["Actor"] --> R["Result"]
  I2["Input B"] --> A
  C([Condition]) -. applies .- A
  RS["decisionRule\nbranch 1 ...\nbranch 2 ..."] -.- A
\end{lstlisting}

\begin{ddtacore}[title={Semantic determinism e omissione controllata}]
A parita' di BA e projection profile, la decisione di includere o omettere \texttt{operatorStructure} deve essere deterministica. Se inclusa, branch/label/assignment derivano dalla BA; se omessa, la BA resta invariata e la view dichiara l'omissione tramite il proprio coverage contract. Nessun modello ``collapsed'' viene creato.
\end{ddtacore}
%</PAGE-CONTENT:75>

\clearpage

%<PAGE-DESC:76>12 - decisionRule: esempio grafico DermaTriage candidate
%<PAGE-CONTENT:76>
\PageHashFooter{\PageMDfive{76}}
\subsection{Esempio grafico corrispondente}
La figura usa soltanto il \texttt{produce} necessario al pressure test DermaTriage. Non espande Stage 1--4 o altre proposition. Il box grigio tratteggiato e' struttura interna renderizzata; una view che non necessita della decision logic puo' ometterlo completamente.

\vspace{3mm}
\begin{center}
\begin{tikzpicture}[
  actor/.style={draw=DDTABlue,fill=DDTALightBlue,rounded corners=3pt,thick,align=center,minimum height=17mm,text width=38mm,inner sep=5pt,font=\sffamily\fontsize{6.8}{8.0}\selectfont},
  contentdoc/.style={draw=DDTABlue,fill=white,thick,align=center,minimum height=18mm,text width=34mm,inner sep=5pt,font=\sffamily\fontsize{6.6}{7.7}\selectfont,
    path picture={\draw[DDTABlue,thick] ([xshift=-5mm]path picture bounding box.north east) -- ([xshift=-5mm,yshift=-5mm]path picture bounding box.north east) -- ([yshift=-5mm]path picture bounding box.north east); \draw[DDTABlue,thick] ([xshift=-5mm]path picture bounding box.north east) -- ([yshift=-5mm]path picture bounding box.north east);}},
  flow/.style={draw=DDTABlue,very thick,-{Stealth[length=2.8mm,width=1.9mm]}},
  conditiontag/.style={ellipse,draw=DDTAAmber,fill=DDTALightAmber,thick,align=center,minimum width=38mm,minimum height=10mm,inner sep=3pt,font=\sffamily\fontsize{6.7}{7.8}\selectfont},
  rules/.style={draw=DDTAGray,fill=DDTALightGray,dashed,rounded corners=2pt,thick,align=left,text width=82mm,inner sep=7pt,font=\ttfamily\fontsize{6.4}{7.5}\selectfont},
  qualifier/.style={draw=DDTAGray,thick,dashed}
]
\node[contentdoc] (u) at (-5.1,1.1) {\textbf{CaseUrgency}\\{\scriptsize\texttt{input}}};
\node[contentdoc] (c) at (-5.1,-1.1) {\textbf{TriageConfidence}\\{\scriptsize\texttt{input}}};
\node[actor] (a) at (0,0) {\textbf{OperationalPriority}\\\textbf{Derivation}\\{\scriptsize\texttt{actor - candidate}}};
\node[contentdoc] (r) at (5.2,0) {\textbf{OperationalTriagePriority}\\{\scriptsize\texttt{result}}};
\node[conditiontag] (g) at (0,2.4) {\texttt{LesionImage available}};
\node[rules] (rs) at (0,-3.7) {\textbf{decisionRule}\\[1mm]
HIGH + confidence $>$ 0.85 $\rightarrow$ P1\\
HIGH, other applicable cases $\rightarrow$ P2\\
MEDIUM $\rightarrow$ P3\\
LOW $\rightarrow$ P4};
\draw[flow] (u.east) -- (a.west);
\draw[flow] (c.east) -- (a.west);
\draw[flow] (a.east) -- (r.west);
\draw[qualifier] (g.south) -- (a.north);
\draw[qualifier] (a.south) -- (rs.north);
\end{tikzpicture}
\end{center}

\begin{ddtacheck}[title={Lettura della candidate view}]
Blu = produzione; ambra = applicability; grigio tratteggiato = result selection interna. Il collegamento grigio non e' dataflow. La figura non promuove la rilocazione: visualizza il \texttt{decisionRule} annidato candidate da sottoporre a regression prima di modificare l'authority BA.
\end{ddtacheck}
%</PAGE-CONTENT:76>

\clearpage

%<PAGE-DESC:77>Appendice temporanea - scopo, lineage e confine del rebuild
%<PAGE-CONTENT:77>
\PageHashFooter{\PageMDfive{77}}
\appendix
\section{Materiale temporaneo di rebuild}
\TemporaryRebuildBanner
\subsection{Scopo, lineage e confine della ricostruzione}

Questa guida non nasce per dichiarare che la revisione piu' recente sia automaticamente la migliore. Nasce per rendere leggibile e applicabile il patrimonio BA gia' costruito, riesaminandolo in modo uniforme mentre viene applicato alla documentazione DermaTriage corrente.

\begin{ddtacore}[title={Regola di costruzione cumulativa}]
La nuova struttura editoriale puo' essere diversa da R3, Core R42 o Complete R42. Il significato precedente, pero', non puo' scomparire per effetto della riscrittura. Ogni distinzione precedente deve essere preservata, raffinata, rilocalizzata, ritestata, respinta con motivazione oppure lasciata esplicitamente aperta.
\end{ddtacore}

\subsection{Tre livelli da non confondere}

\begin{tabularx}{\textwidth}{L{38mm}L{39mm}Y}
\toprule
\textbf{Livello} & \textbf{Esempio} & \textbf{Effetto} \\
\midrule
Authority corrente & BA0--BA5 + Operational Guide R3 & Governa il metodo corrente finche' non esiste promotion esplicita. \\
Research evidence & R42, operator review, ledger, pressure, candidate, vecchi appunti & Puo' motivare chiarimenti o cambi, ma non modifica da solo il metodo. \\
Rebuild candidate & questa guida & Superficie cumulativa di review; diventa authority solo con checkpoint successivo. \\
\bottomrule
\end{tabularx}

\begin{ddtaopen}[title={Nessun target di cardinalita'}]
Il fatto che la BA corrente abbia quattordici operatori top-level non crea un obiettivo di quattordici operatori. Anche il precedente snapshot a tredici elementi resta genealogia utile, non un target da ripristinare. Il numero finale deve essere un risultato della necessita' semantica e della regression, non una quota.
\end{ddtaopen}
%</PAGE-CONTENT:77>

\clearpage

%<PAGE-DESC:78>Appendice temporanea - corpus e stati epistemici
%<PAGE-CONTENT:78>
\PageHashFooter{\PageMDfive{78}}
\TemporaryRebuildBanner
\subsection{Corpus di ricostruzione e stati epistemici}

Il rebuild usa come input controllato il registro {\scriptsize\nolinkurl{DDTA_R25_BASE_ANALYSIS_REBUILD_SOURCE_REGISTER_R1.md}}. Prima di chiudere un costrutto vanno ricercati non solo il suo nome corrente, ma anche alias storici, pressure collegate, candidate construct, alternative rifiutate e counterexample.

\begin{ddtaprinciple}[title={Recency non equivale ad authority}]
R42 Core/Complete e' una superficie di lettura aggiornata e molto utile per il rebuild. R3 e BA0--BA5 restano tuttavia authority corrente. La guida nuova puo' usare R42 come base esplicativa senza promuovere automaticamente le sue candidate o chiudere le sue open pressure.
\end{ddtaprinciple}

\subsection{Etichette epistemiche minime}

\begin{tabularx}{\textwidth}{L{49mm}L{34mm}Y}
\toprule
\textbf{Etichetta} & \textbf{Uso} & \textbf{Disciplina} \\
\midrule
\CurrentTag & Semantica corrente preservata & Deve restare compatibile con l'authority o dichiararne esplicitamente il delta. \\
\CandidateTag & Costrutto/evoluzione con evidence positiva ma non ammessa & Non usare come vocabulary corrente senza admission esplicita. \\
\OpenTag & Problema reale non chiuso & Mantenere il gap visibile; non completare per convenienza. \\
\RejectedTag & Ipotesi testata e respinta nello scope disponibile & Conservare ragione, test e condizioni di eventuale reopen. \\
\RebuildTag & Regola editoriale del rebuild & Governa come descriviamo e confrontiamo i costrutti durante questa ricostruzione. \\
\bottomrule
\end{tabularx}

\begin{ddtacheck}[title={Regola anti-perdita}]
Una guida piu' corta non e' una guida migliore se per ottenere la sintesi elimina una distinzione ancora utile. Il rebuild puo' comprimere la prosa soltanto dopo aver dimostrato che la distinzione resta ricostruibile e tracciabile.
\end{ddtacheck}
%</PAGE-CONTENT:78>

\clearpage

%<PAGE-DESC:79>Appendice temporanea - Construct Description Contract e principio espositivo
%<PAGE-CONTENT:79>
\PageHashFooter{\PageMDfive{79}}
\TemporaryRebuildBanner
\subsection{Construct Description Contract}

\RebuildTag

Prima di riscrivere il primo costrutto BA, fissiamo il modo in cui \emph{tutti} i costrutti verranno descritti. Questo contratto serve a impedire che la guida cambi struttura in funzione del caso del momento e rende confrontabili operatori current, condition form, strutture candidate, pressure e ipotesi rifiutate.

\begin{ddtacore}[title={Invariante editoriale}]
Ogni costrutto usa la stessa sequenza descrittiva. Un campo puo' essere \texttt{NOT APPLICABLE}, \texttt{NOT FROZEN} o \texttt{OPEN}; non puo' essere omesso silenziosamente. Se un costrutto futuro dimostra che il contratto non basta, si modifica prima il contratto e poi si riallineano i costrutti gia' trattati.
\end{ddtacore}

\begin{ddtaprinciple}[title={Controllo rigoroso, esposizione naturale}]
Il CDC e' un contratto di \emph{copertura}, non un template rigido di impaginazione. I campi richiesti devono essere tutti trattati e ricostruibili durante la review, ma la guida definitiva puo' combinarli in paragrafi e sezioni piu' naturali quando questo non nasconde informazioni, stato epistemico o confini semantici. Struttura di controllo rigorosa non significa esposizione meccanica.
\end{ddtaprinciple}
%</PAGE-CONTENT:79>

\clearpage

%<PAGE-DESC:80>Appendice temporanea - CDC direzione della spiegazione e revisione
%<PAGE-CONTENT:80>
\PageHashFooter{\PageMDfive{80}}
\TemporaryRebuildBanner
\subsection{Direzione obbligatoria della spiegazione}

\begin{lstlisting}
human semantic idea
    -> BA term / construct class
    -> formal shape
    -> boundaries and discriminators
    -> examples
    -> evidence and genealogy
    -> open questions
    -> current rebuild disposition
\end{lstlisting}

\begin{ddtaanti}[title={Da evitare}]
\begin{itemize}
  \item iniziare dalla firma e lasciare al lettore il compito di indovinarne il significato;
  \item usare l'esempio DermaTriage come definizione implicita del costrutto;
  \item aggiungere campi solo a un costrutto perche' il suo caso concreto ne ha bisogno;
  \item completare la firma di una candidate non ammessa per rendere il template piu' ordinato;
  \item nascondere rejected hypothesis o pressure per ottenere una sezione apparentemente ``chiusa''.
\end{itemize}
\end{ddtaanti}

\subsection{Revisione del contratto}
Il contratto parte come \texttt{CDC-R1}. Ogni sua modifica futura deve avere un counterexample, una motivazione generalizzabile e una regression sui costrutti gia' descritti.
%</PAGE-CONTENT:80>

\clearpage

%<PAGE-DESC:81>Appendice temporanea - CDC campi 1-10
%<PAGE-CONTENT:81>
\PageHashFooter{\PageMDfive{81}}
\TemporaryRebuildBanner
\subsection{Construct Description Contract - template, parte 1}

\begin{longtable}{L{10mm}L{47mm}L{84mm}}
\toprule
\textbf{N.} & \textbf{Campo} & \textbf{Contenuto richiesto} \\
\midrule\endhead
1 & Canonical name & Nome usato dal rebuild. Alias storici vanno nel campo genealogy, non come nomi paralleli current. \\
2 & Construct class & Top-level operator, condition form, local reusable structure, candidate, open pressure, rejected hypothesis o altro meccanismo BA dichiarato. \\
3 & Epistemic status & Stato esplicito: current/preserved, candidate/not admitted, open, rejected oppure altro stato definito. \\
4 & Owning BA layer / contract & Dove il significato appartiene: BA0, BA1, BA2, BA3, BA4, BA5 oppure confine cross-layer esplicito. \\
5 & Human definition & Spiegazione in linguaggio naturale di che cosa rappresenta il costrutto. \\
6 & Smallest semantic fact & Il nucleo minimo che si perde eliminando il costrutto. \\
7 & When to use & Condizioni positive che giustificano l'uso. Devono riferirsi a significato governato, non a parole chiave. \\
8 & When not to use & Casi in cui il costrutto sembra plausibile ma imporrebbe semantica sbagliata o eccessiva. \\
9 & Canonical syntax / shape & Forma corrente se ammessa. Per candidate non chiuse: \texttt{NOT FROZEN} salvo working shape esplicitamente marcata. \\
10 & Roles, types, cardinalities & Ruoli e vincoli formali ammessi. Nessun role viene inventato per l'esempio corrente. \\
\bottomrule
\end{longtable}

\begin{ddtaprinciple}[title={Semantica prima della sintassi}]
I campi 5--8 devono poter essere compresi anche da chi non conosce ancora la sintassi BA. I campi 9--10 formalizzano il significato; non lo sostituiscono.
\end{ddtaprinciple}
%</PAGE-CONTENT:81>

\clearpage

%<PAGE-DESC:82>Appendice temporanea - CDC campi 11-20
%<PAGE-CONTENT:82>
\PageHashFooter{\PageMDfive{82}}
\TemporaryRebuildBanner
\subsection{Construct Description Contract - template, parte 2}

\begin{longtable}{L{10mm}L{47mm}L{84mm}}
\toprule
\textbf{N.} & \textbf{Campo} & \textbf{Contenuto richiesto} \\
\midrule\endhead
11 & Semantic invariants & Cio' che deve restare vero in ogni uso valido del costrutto. \\
12 & Forbidden inferences & Fatti che la presenza del costrutto non autorizza a dedurre. \\
13 & Closest constructs / discriminators & Confini pairwise con i costrutti piu' facili da confondere. \\
14 & Minimal generic positive example & Esempio didattico minimo, separato dall'evidence di progetto. \\
15 & DermaTriage application example & Uso reale sulla documentazione corrente, quando disponibile e supportato. \\
16 & Negative / boundary example & Caso in cui il costrutto va rifiutato o usato con limite preciso. \\
17 & Genealogy and research evidence & Authority, candidate guide, operator review, PR/CC/CL/CMD/OBS e precedenti rilevanti. \\
18 & Open pressures / candidate deltas & Questioni ancora vive che non devono essere nascoste dalla descrizione principale. \\
19 & Review / acceptance questions & Domande ripetibili che permettono a un autore di verificare l'uso. \\
20 & Current rebuild disposition & Cosa decide questa revisione: preserve, refine, retest, reject, hold, candidate admission later, ecc. \\
\bottomrule
\end{longtable}

\begin{ddtacheck}[title={Change protocol del CDC}]
Se un costrutto non entra nel template, non si crea una eccezione locale. Si registra il counterexample, si verifica se il problema e' generale, si produce una revisione \texttt{CDC-R2} e si riesaminano retroattivamente i costrutti gia' descritti. Nessuna modifica implicita del formato e' ammessa.
\end{ddtacheck}

\begin{ddtaexample}[title={Candidate con firma non chiusa}]
Per \texttt{provideService} il template resta completo, ma i campi ``Canonical syntax'' e ``Roles, types, cardinalities'' possono dichiarare \texttt{NOT FROZEN}. E' preferibile una lacuna esplicita a una firma inventata solo per uniformare l'impaginazione.
\end{ddtaexample}
%</PAGE-CONTENT:82>

\clearpage

%<PAGE-DESC:83>Appendice temporanea - roadmap corrente del rebuild R7
%<PAGE-CONTENT:83>
\PageHashFooter{\PageMDfive{83}}
\TemporaryRebuildBanner
\subsection{Roadmap temporanea del rebuild}
\begin{enumerate}
  \item usare \texttt{CDC-R1} come controllo di copertura, senza trasformarlo in una scheda editoriale rigida;
  \item assumere le Sezioni 1--12 come base espositiva corrente della Rebuild R7;
  \item preservare governed semantic fact, \BAR{}, \BAP{}, identity criterion, stopping rule e structural-feedback boundary come fondazioni gia' riesaminate;
  \item considerare \texttt{transfer}, \texttt{produce}, \texttt{realize} e \texttt{constrain} i primi quattro operatori completamente ricostruiti, \texttt{condition} il primo scoped modifier completamente ricostruito e \texttt{decisionRule} la prima \texttt{operatorStructure} candidate riesaminata sotto lo stesso \texttt{CDC-R1};
  \item mantenere come placeholder le discriminazioni pairwise dei quattro operatori dove questa scelta resta deliberata; per \texttt{condition} e la candidate annidata \texttt{decisionRule}, conservare invece i discriminatori necessari fra applicability, result selection, allowed-domain constraint, dependency e identity pressure;
  \item preservare in \texttt{constrain} la distinzione fra constraint nominato/riusabile e restriction locale, e nella projection distinguere il vincolo sul content dal vincolo sul transfer behavior;
  \item continuare la Base Analysis di \texttt{MR-01} dalla documentazione DermaTriage corrente consolidata in R25; usare \texttt{FR-MR01-04-01} come pressure case per la rilocazione di \texttt{decisionRule} senza trattare la candidate come authority prima della regression;
  \item per il prossimo costrutto effettivamente necessario, partire dalle guide precedenti, ricostruire prima il metodo generale e poi verificarlo con evidence corrente DermaTriage;
  \item mantenere prompt LLM e wrapper di authoring separati dalla semantica BA: possono evolvere senza creare nuovi campi o operatori;
  \item solo dopo BA accettata per uno scope, produrre e verificare le projection downstream previste.
\end{enumerate}

\begin{ddtaopen}[title={Stop point di questa R7}]
La Rebuild R7 contiene la ricostruzione completa dei quattro operatori \texttt{transfer}, \texttt{produce}, \texttt{realize}, \texttt{constrain} e dello scoped modifier \texttt{condition}. La Sezione 12 aggiunge una candidate strutturale: preservare il mapping di risultato come \texttt{operatorStructure.decisionRule} dell'owner \texttt{produce}, invece di duplicare actor/input/result in una seconda top-level proposition \texttt{decisionRule}. Il delta non e' ammesso: deve ancora superare la regression sul caso Facial Access, il pressure test DermaTriage e un checkpoint esplicito.
\end{ddtaopen}
%</PAGE-CONTENT:83>

\clearpage

%<PAGE-DESC:84>Indice di integrita' delle pagine
%<PAGE-CONTENT:84>
\PageHashFooter{\PageMDfive{84}}
\section*{Indice di integrita' delle pagine}
{\footnotesize Ogni pagina possiede un MD5 del proprio contenuto canonico; il valore stampato e' escluso dal payload e l'hash della pagina indice usa \texttt{<SELF>} per evitare dipendenza circolare.}\par
\vspace{1mm}
\begingroup
\tiny
\renewcommand{\arraystretch}{0.62}
\setlength{\tabcolsep}{1.5pt}
\renewcommand{\IntegrityRow}[3]{#1 & #2 & {\tiny\ttfamily #3} \\}
\begin{longtable}{L{7mm}L{93mm}L{44mm}}
\toprule
\textbf{Pag.} & \textbf{Contenuto} & \textbf{MD5 contenuto} \\
\midrule
%<PAGE-INTEGRITY-ROWS>
\IntegrityRow{1}{Frontespizio della Base Analysis Guide Rebuild R7}{\PageMDfive{1}}
\IntegrityRow{2}{1 - Che cos'e' la Base Analysis}{\PageMDfive{2}}
\IntegrityRow{3}{1 - Perche' esiste e principi di base}{\PageMDfive{3}}
\IntegrityRow{4}{1 - Direzione dell'authority e feedback}{\PageMDfive{4}}
\IntegrityRow{5}{1 - Minimum sufficient BA, NOT SPECIFIED e gap}{\PageMDfive{5}}
\IntegrityRow{6}{2 - BA View / Projection Contract e source types}{\PageMDfive{6}}
\IntegrityRow{7}{2 - Mermaid e campi minimi della view}{\PageMDfive{7}}
\IntegrityRow{8}{2 - Coverage mode e famiglie di viste}{\PageMDfive{8}}
\IntegrityRow{9}{2 - View gap e mapping discipline}{\PageMDfive{9}}
\IntegrityRow{10}{3 - Dal testo documentale al significato governato}{\PageMDfive{10}}
\IntegrityRow{11}{3 - Governed semantic fact}{\PageMDfive{11}}
\IntegrityRow{12}{3 - Segmentazione, ambiguita e livello documentale}{\PageMDfive{12}}
\IntegrityRow{13}{4 - Elementi fondamentali della Base Analysis}{\PageMDfive{13}}
\IntegrityRow{14}{4 - BAReferent: identita analitica riusabile}{\PageMDfive{14}}
\IntegrityRow{15}{4 - BAReferent: esempi DermaTriage e confini}{\PageMDfive{15}}
\IntegrityRow{16}{4 - BAProposition: asserzione analitica identificabile}{\PageMDfive{16}}
\IntegrityRow{17}{4 - BAReferent e BAProposition insieme}{\PageMDfive{17}}
\IntegrityRow{18}{5 - Granularita, minimalita e costruzione della BA}{\PageMDfive{18}}
\IntegrityRow{19}{6 - Indipendenza tra struttura documentale e struttura BA}{\PageMDfive{19}}
\IntegrityRow{20}{6 - Structural feedback e no authority inversion}{\PageMDfive{20}}
\IntegrityRow{21}{6 - DermaTriage DEC02 DEC04 come pressure case}{\PageMDfive{21}}
\IntegrityRow{22}{6 - Downstream Utility Test e riallineamento BA}{\PageMDfive{22}}
\IntegrityRow{23}{7 - transfer: definizione, smallest fact e criteri d'uso}{\PageMDfive{23}}
\IntegrityRow{24}{7 - transfer: firma canonica, ruoli e cardinalita}{\PageMDfive{24}}
\IntegrityRow{25}{7 - transfer: identity continuity, invarianti e forbidden inference}{\PageMDfive{25}}
\IntegrityRow{26}{7 - transfer: domande di identificazione e procedura di estrazione}{\PageMDfive{26}}
\IntegrityRow{27}{7 - transfer: esempio generico e worked example DermaTriage corrente}{\PageMDfive{27}}
\IntegrityRow{28}{7 - transfer: prompt LLM per l'identificazione controllata}{\PageMDfive{28}}
\IntegrityRow{29}{7 - transfer: output BA atteso e discriminazioni differite}{\PageMDfive{29}}
\IntegrityRow{30}{7 - transfer: rappresentazione Mermaid}{\PageMDfive{30}}
\IntegrityRow{31}{7 - transfer: esempio grafico Mermaid}{\PageMDfive{31}}
\IntegrityRow{32}{8 - produce: definizione, smallest fact e criteri d'uso}{\PageMDfive{32}}
\IntegrityRow{33}{8 - produce: firma canonica, ruoli e cardinalita}{\PageMDfive{33}}
\IntegrityRow{34}{8 - produce: identity continuity, invarianti e forbidden inference}{\PageMDfive{34}}
\IntegrityRow{35}{8 - produce: domande di identificazione e procedura di estrazione}{\PageMDfive{35}}
\IntegrityRow{36}{8 - produce: esempio generico e worked example DermaTriage corrente}{\PageMDfive{36}}
\IntegrityRow{37}{8 - produce: prompt LLM per l'identificazione controllata}{\PageMDfive{37}}
\IntegrityRow{38}{8 - produce: output BA atteso e discriminazioni differite}{\PageMDfive{38}}
\IntegrityRow{39}{8 - produce: rappresentazione Mermaid}{\PageMDfive{39}}
\IntegrityRow{40}{8 - produce: esempio grafico Mermaid}{\PageMDfive{40}}
\IntegrityRow{41}{9 - realize: definizione, smallest fact e criteri d'uso}{\PageMDfive{41}}
\IntegrityRow{42}{9 - realize: firma canonica, ruoli e cardinalita}{\PageMDfive{42}}
\IntegrityRow{43}{9 - realize: identity continuity, invarianti e forbidden inference}{\PageMDfive{43}}
\IntegrityRow{44}{9 - realize: domande di identificazione e procedura di estrazione}{\PageMDfive{44}}
\IntegrityRow{45}{9 - realize: esempio generico e worked example DermaTriage corrente}{\PageMDfive{45}}
\IntegrityRow{46}{9 - realize: prompt LLM per l'identificazione controllata}{\PageMDfive{46}}
\IntegrityRow{47}{9 - realize: output BA atteso e discriminazioni differite}{\PageMDfive{47}}
\IntegrityRow{48}{9 - realize: rappresentazione Mermaid}{\PageMDfive{48}}
\IntegrityRow{49}{9 - realize: esempio grafico Mermaid}{\PageMDfive{49}}
\IntegrityRow{50}{10 - constrain: definizione, smallest fact e criteri d'uso}{\PageMDfive{50}}
\IntegrityRow{51}{10 - constrain: firma canonica e forme di constraintValue}{\PageMDfive{51}}
\IntegrityRow{52}{10 - constrain: identity continuity, invarianti e forbidden inference}{\PageMDfive{52}}
\IntegrityRow{53}{10 - constrain: domande di identificazione e procedura di estrazione}{\PageMDfive{53}}
\IntegrityRow{54}{10 - constrain: esempio generico e worked example DermaTriage corrente}{\PageMDfive{54}}
\IntegrityRow{55}{10 - constrain: prompt LLM per l'identificazione controllata}{\PageMDfive{55}}
\IntegrityRow{56}{10 - constrain: output BA atteso e discriminazioni differite}{\PageMDfive{56}}
\IntegrityRow{57}{10 - constrain: rappresentazione Mermaid}{\PageMDfive{57}}
\IntegrityRow{58}{10 - constrain: esempio grafico Mermaid composto con transfer}{\PageMDfive{58}}
\IntegrityRow{59}{11 - condition: definizione, smallest fact e criteri d'uso}{\PageMDfive{59}}
\IntegrityRow{60}{11 - condition: forma canonica, scope locale e cardinalita}{\PageMDfive{60}}
\IntegrityRow{61}{11 - condition: identity boundary, invarianti e forbidden inference}{\PageMDfive{61}}
\IntegrityRow{62}{11 - condition: domande di identificazione e procedura di estrazione}{\PageMDfive{62}}
\IntegrityRow{63}{11 - condition: esempio generico e worked example DermaTriage corrente}{\PageMDfive{63}}
\IntegrityRow{64}{11 - condition: prompt LLM per l'identificazione controllata}{\PageMDfive{64}}
\IntegrityRow{65}{11 - condition: output BA atteso e discriminazioni}{\PageMDfive{65}}
\IntegrityRow{66}{11 - condition: rappresentazione Mermaid e semantic determinism}{\PageMDfive{66}}
\IntegrityRow{67}{11 - condition: esempio grafico Mermaid composto con produce}{\PageMDfive{67}}
\IntegrityRow{68}{12 - decisionRule: definizione, smallest fact e criteri d'uso}{\PageMDfive{68}}
\IntegrityRow{69}{12 - decisionRule: forma canonica candidate e cardinalita}{\PageMDfive{69}}
\IntegrityRow{70}{12 - decisionRule: identity boundary, invarianti e forbidden inference}{\PageMDfive{70}}
\IntegrityRow{71}{12 - decisionRule: domande di identificazione e procedura di estrazione}{\PageMDfive{71}}
\IntegrityRow{72}{12 - decisionRule: esempio generico e pressure case DermaTriage}{\PageMDfive{72}}
\IntegrityRow{73}{12 - decisionRule: prompt LLM per l'identificazione controllata}{\PageMDfive{73}}
\IntegrityRow{74}{12 - decisionRule: output candidate e discriminazioni}{\PageMDfive{74}}
\IntegrityRow{75}{12 - decisionRule: rappresentazione Mermaid e visibility profile}{\PageMDfive{75}}
\IntegrityRow{76}{12 - decisionRule: esempio grafico DermaTriage candidate}{\PageMDfive{76}}
\IntegrityRow{77}{Appendice temporanea - scopo, lineage e confine del rebuild}{\PageMDfive{77}}
\IntegrityRow{78}{Appendice temporanea - corpus e stati epistemici}{\PageMDfive{78}}
\IntegrityRow{79}{Appendice temporanea - Construct Description Contract e principio espositivo}{\PageMDfive{79}}
\IntegrityRow{80}{Appendice temporanea - CDC direzione della spiegazione e revisione}{\PageMDfive{80}}
\IntegrityRow{81}{Appendice temporanea - CDC campi 1-10}{\PageMDfive{81}}
\IntegrityRow{82}{Appendice temporanea - CDC campi 11-20}{\PageMDfive{82}}
\IntegrityRow{83}{Appendice temporanea - roadmap corrente del rebuild R7}{\PageMDfive{83}}
\IntegrityRow{84}{Indice di integrita' delle pagine}{\PageMDfive{84}}
\IntegrityRow{85}{Appendice temporanea - laboratorio prove grafiche Mermaid da cancellare}{\PageMDfive{85}}
%</PAGE-INTEGRITY-ROWS>
\bottomrule
\end{longtable}
\endgroup
\vspace{1mm}
{\footnotesize La pagina successiva resta il laboratorio grafico temporaneo da cancellare; le altre appendici di rebuild restano deliberate fino al consolidamento.}
%</PAGE-CONTENT:84>


\clearpage

%<PAGE-DESC:85>Appendice temporanea - laboratorio prove grafiche Mermaid da cancellare
%<PAGE-CONTENT:85>
\PageHashFooter{\PageMDfive{85}}
\section*{PROVA GRAFICA TEMPORANEA - DA CANCELLARE}

\begin{ddtaopen}[title={Scopo della pagina}]
Questa pagina e' un laboratorio grafico temporaneo e non stabilisce da sola le projection definitive. Deve essere eliminata quando le convenzioni saranno stabilizzate. La prova corrente conserva la vista composta \texttt{produce + transfer} nella quale \texttt{ClinicalDescription} resta un dato passivo: e' result della produzione e content del transfer, ma non genera frecce verso la destination. Le prove di operatori successivi possono essere aggiunte qui senza trasformare questa appendice in metodologia normativa.
\end{ddtaopen}

\subsection*{Prova A - content visivamente tra source e transfer}
La disposizione cercata mette \texttt{ClinicalDescription} fisicamente fra source e \texttt{transfer}. Il percorso attivo resta \texttt{source -> transfer -> destination}; la linea tratteggiata \texttt{content} e' una associazione di role e non un flusso.

\begin{lstlisting}
flowchart LR
  S["ClinicalDescriptionGeneration<br/>actor / source"]
  C@{ shape: doc, label: "ClinicalDescription<br/>result / content" }
  T(("transfer"))
  D["HistoricalCaseRetrieval<br/>destination"]
  B["ClinicalDescriptionTransfer<br/>behavior"]

  S -->|"produce"| C
  S ==> T ==> D
  C -.-|"content"| T
  B -.-|"behavior"| T

  classDef endpoint fill:#EEF4F9,stroke:#275D86,color:#17324D
  classDef content fill:#FFFFFF,stroke:#275D86,color:#17324D
  classDef transfer fill:#DCEAF5,stroke:#275D86,color:#17324D
  classDef behavior fill:#FFF4E3,stroke:#A96A16,color:#17324D
  class S,D endpoint
  class C content
  class T transfer
  class B behavior
\end{lstlisting}

\vspace{2mm}
\begin{center}
\begin{tikzpicture}[
  endpoint/.style={draw=DDTABlue,fill=DDTALightBlue,rounded corners=3pt,thick,align=center,minimum height=14mm,text width=31mm,inner sep=4pt,font=\sffamily\fontsize{6.7}{7.7}\selectfont},
  contentdoc/.style={draw=DDTABlue,fill=white,thick,align=center,minimum height=18mm,text width=28mm,inner sep=5pt,font=\sffamily\fontsize{6.7}{7.7}\selectfont,
    path picture={\draw[DDTABlue,thick] ([xshift=-5mm]path picture bounding box.north east) -- ([xshift=-5mm,yshift=-5mm]path picture bounding box.north east) -- ([yshift=-5mm]path picture bounding box.north east); \draw[DDTABlue,thick] ([xshift=-5mm]path picture bounding box.north east) -- ([yshift=-5mm]path picture bounding box.north east);}},
  transfernode/.style={circle,draw=DDTABlue,fill=DDTALightBlue,thick,align=center,minimum size=17mm,inner sep=2pt,font=\sffamily\bfseries\fontsize{7}{8}\selectfont},
  behavior/.style={draw=DDTAAmber,fill=DDTALightAmber,rounded corners=3pt,thick,align=center,minimum height=12mm,text width=36mm,inner sep=4pt,font=\sffamily\fontsize{6.7}{7.7}\selectfont},
  produceflow/.style={draw=DDTABlue,thick,-{Stealth[length=2.6mm,width=1.8mm]}},
  transferflow/.style={draw=DDTABlue,very thick,-{Stealth[length=3mm,width=2mm]}},
  passive/.style={draw=DDTAGray,thick,dashed}
]
\node[endpoint] (s) at (-5.0,0) {\textbf{ClinicalDescription}\\\textbf{Generation}\\{\scriptsize\texttt{actor / source}}};
\node[contentdoc] (c) at (-1.9,1.05) {\textbf{ClinicalDescription}\\{\scriptsize\texttt{result / content}}};
\node[transfernode] (t) at (1.2,0) {transfer};
\node[endpoint] (d) at (4.9,0) {\textbf{HistoricalCaseRetrieval}\\{\scriptsize\texttt{destination}}};
\node[behavior] (b) at (0.8,-2.15) {\textbf{ClinicalDescriptionTransfer}\\{\scriptsize\texttt{behavior}}};

\draw[produceflow] (s.north east) -- (c.west);\node[font=\sffamily\scriptsize] at (-3.55,1.05) {produce};
\draw[transferflow] (s.east) -- (t.west);
\draw[transferflow] (t.east) -- (d.west);
\draw[passive] (c.south east) -- node[right,font=\sffamily\scriptsize]{content} (t.north west);
\draw[passive] (b.north east) -- node[right,font=\sffamily\scriptsize]{behavior} (t.south west);
\end{tikzpicture}
\end{center}


%</PAGE-CONTENT:85>

\end{document}
